\documentclass[a4paper,11pt]{article}

\usepackage[ngerman]{babel}
\usepackage[T1]{fontenc}
\usepackage[utf8]{inputenc}
\usepackage[a4paper,margin=1.8cm]{geometry}
\usepackage{amsmath,amssymb}
\usepackage{enumitem}
\usepackage[hidelinks]{hyperref}

\setlist[itemize]{leftmargin=1.6em, itemsep=0.2em, topsep=0.35em}
\setlist[enumerate]{leftmargin=1.8em, itemsep=0.2em, topsep=0.35em}
\setlength{\parindent}{0pt}
\setlength{\parskip}{0.45em}

\title{Stochastik Cheatsheet}
\author{laufende Sammlung fuer die Klausur}
\date{\today}

\begin{document}

\maketitle
\tableofcontents
\bigskip

\section{Kapitel 1.1: Zufallsexperimente und Wahrscheinlichkeitsraeume}

\subsection{Grundnotation}
\begin{itemize}
  \item $\Omega$: Ergebnismenge, also die Menge aller moeglichen Ergebnisse.
  \item $\omega \in \Omega$: ein einzelnes Ergebnis.
  \item $A, B, C \subseteq \Omega$: Ereignisse.
  \item $\emptyset$: unmoegliches Ereignis.
  \item $\Omega$: sicheres Ereignis.
  \item $\{\omega\}$: Elementarereignis zum Ergebnis $\omega$.
  \item $\overline{A} = \Omega \setminus A$: Gegenereignis zu $A$.
  \item $A \cup B$: ``$A$ oder $B$''.
  \item $A \cap B$: ``$A$ und $B$''.
  \item $A \setminus B$: ``$A$ ohne $B$''.
\end{itemize}

\textbf{Merke:}
Ein Ergebnis ist ein einzelner Ausgang. Ein Ereignis ist eine Menge von Ergebnissen. Ein Elementarereignis ist ein Ereignis mit genau einem Ergebnis.

\subsection{Zufallsexperiment}
Ein Vorgang ist ein Zufallsexperiment, wenn:
\begin{itemize}
  \item er unter gleichen Bedingungen wiederholbar ist,
  \item er ein Ergebnis liefert,
  \item die moeglichen Ergebnisse vorher feststehen,
  \item das konkrete Ergebnis nicht sicher vorhersagbar ist.
\end{itemize}

\subsection{Ereignisse sauber notieren}
Beispiel beim Wuerfeln mit $\Omega = \{1,2,3,4,5,6\}$:
\begin{itemize}
  \item gerade Zahl: $A = \{2,4,6\}$
  \item hoechstens $4$: $B = \{1,2,3,4\}$
  \item gerade und hoechstens $4$: $A \cap B = \{2,4\}$
  \item gerade oder hoechstens $4$: $A \cup B = \{1,2,3,4,6\}$
\end{itemize}

\textbf{Merke:}
Text immer zuerst in Mengen uebersetzen, dann rechnen.

\subsection{Wahrscheinlichkeitsmass}
Ein Wahrscheinlichkeitsmass ist eine Abbildung
\[
P : \mathcal{P}(\Omega) \to [0,1]
\]
mit:
\begin{align*}
P(\Omega) &= 1, \\
P\left(\bigcup_{i=1}^{\infty} A_i\right) &= \sum_{i=1}^{\infty} P(A_i)
\quad \text{fuer paarweise disjunkte } A_1, A_2, \dots
\end{align*}

\textbf{Wichtige Rechenregeln}
\begin{align*}
P(\emptyset) &= 0 \\
P(\overline{A}) &= 1 - P(A) \\
A \subseteq B &\Rightarrow P(A) \leq P(B) \\
P(A \setminus B) &= P(A) - P(A \cap B) \\
P(A \cup B) &= P(A) + P(B) - P(A \cap B)
\end{align*}

Falls $A \cap B = \emptyset$, also $A$ und $B$ disjunkt sind:
\[
P(A \cup B) = P(A) + P(B)
\]

\textbf{Merke:}
Bei ``oder'' fast immer zuerst an
\[
P(A \cup B) = P(A) + P(B) - P(A \cap B)
\]
denken. Nur wenn die Ereignisse disjunkt sind, faellt der Schnittterm weg.

\subsection{Laplace-Experiment}
Wenn alle Elementarergebnisse gleich wahrscheinlich sind, gilt:
\[
P(A) = \frac{|A|}{|\Omega|}
\]

\textbf{Bedeutung:}
\begin{itemize}
  \item $|A|$: Anzahl guenstiger Ergebnisse
  \item $|\Omega|$: Anzahl aller moeglichen Ergebnisse
\end{itemize}

Beispiel fairer Wuerfel:
\[
\Omega = \{1,2,3,4,5,6\}, \quad A = \{2,4,6\}, \quad P(A) = \frac{3}{6} = \frac{1}{2}
\]

\textbf{Merke:}
Laplace nur benutzen, wenn wirklich alle Ergebnisse gleich wahrscheinlich sind.

\subsection{Wahrscheinlichkeitsfunktion}
Die Wahrscheinlichkeitsfunktion ordnet jedem einzelnen Ergebnis seine Wahrscheinlichkeit zu:
\[
p : \Omega \to [0,1], \quad p(\omega) = P(\{\omega\})
\]

Fuer ein Ereignis $A \subseteq \Omega$ gilt dann:
\[
P(A) = \sum_{\omega \in A} p(\omega)
\]

Damit $p$ eine Wahrscheinlichkeitsfunktion ist, muss gelten:
\begin{align*}
p(\omega) &\geq 0 \quad \text{fuer alle } \omega \in \Omega \\
\sum_{\omega \in \Omega} p(\omega) &= 1
\end{align*}

\textbf{Merke:}
Ist $p$ gegeben, dann werden Ereigniswahrscheinlichkeiten durch Aufsummieren der passenden Einzelwahrscheinlichkeiten berechnet.

\subsection{Standardmuster fuer Aufgaben}
\textbf{1. Pruefen, ob $p$ eine Wahrscheinlichkeitsfunktion ist}
\begin{enumerate}
  \item Zeige $p(\omega) \geq 0$ fuer alle $\omega$.
  \item Zeige $\sum_{\omega \in \Omega} p(\omega) = 1$.
\end{enumerate}

\textbf{2. Konstante $c$ bestimmen}

Wenn $p(i) = c \cdot f(i)$ gegeben ist, dann:
\[
\sum_i p(i) = 1
\quad \Rightarrow \quad
c \sum_i f(i) = 1
\quad \Rightarrow \quad
c = \frac{1}{\sum_i f(i)}
\]

\textbf{3. Wahrscheinlichkeit eines Ereignisses berechnen}

Wenn $A \subseteq \Omega$, dann:
\[
P(A) = \sum_{\omega \in A} p(\omega)
\]

\textbf{4. Formulierungen wie ``mindestens'', ``hoechstens'', ``mehr als''}

Oft ist das Gegenereignis einfacher.

Beispiel:
``10 oder mehr Versuche bis zur ersten 6'' bedeutet:
in den ersten 9 Versuchen faellt keine 6. Also
\[
P(X \geq 10) = \left(\frac{5}{6}\right)^9
\]

\subsection{Notation fuer saubere Klausurloesungen}
Saubere Standardsaetze:
\begin{itemize}
  \item ``Sei $\Omega = \{\dots\}$ die Ergebnismenge.''
  \item ``Sei $A = \{\dots\}$ das Ereignis \dots''
  \item ``Da $A \cap B = \emptyset$, gilt \dots''
  \item ``Mit der Additionsformel folgt \dots''
  \item ``Da es sich um ein Laplace-Experiment handelt, gilt \dots''
  \item ``Fuer die Wahrscheinlichkeitsfunktion muss $\sum_{\omega \in \Omega} p(\omega) = 1$ gelten.''
\end{itemize}

\textbf{Merke:}
Nicht nur Ergebnis hinschreiben, sondern kurz den passenden Ansatz mit notieren. Das gibt eher Punkte.

\subsection{Direkt nuetzlich fuer die Aufgaben des Blatts}
\textbf{Muenzwurf mit Rand}
\[
\Omega = \{K, Z, R\}, \quad p(K) = \frac{1}{2}, \quad p(Z) = \frac{1}{2}, \quad p(R) = 0
\]

\textbf{Wichtig:}
Ein Ergebnis darf in $\Omega$ enthalten sein, auch wenn seine Wahrscheinlichkeit $0$ ist.

\textbf{Konstante Wahrscheinlichkeitsfunktion auf abz\"ahlbar unendlichem $\Omega$}

Unmoeglich:
\begin{itemize}
  \item bei $p(\omega) = c > 0$ waere $\sum p(\omega) = \infty$,
  \item bei $c = 0$ waere $\sum p(\omega) = 0$.
\end{itemize}

Also gibt es auf einer abz\"ahlbar unendlichen Ergebnismenge keine konstante Wahrscheinlichkeitsfunktion.

\textbf{Wuerfeln bis zur ersten 6}

Falls $X$ die Anzahl der Versuche bis zur ersten 6 ist:
\[
P(X = k) = \left(\frac{5}{6}\right)^{k-1} \cdot \frac{1}{6}
\quad \text{fuer } k \in \mathbb{N}
\]

\textbf{Merke:}
Erst $k-1$ mal keine 6, dann einmal 6.

\subsection{Mini-Merkliste}
\begin{itemize}
  \item Erst $\Omega$ aufschreiben, dann Ereignisse definieren, dann rechnen.
  \item ``oder'' ist Vereinigung, ``und'' ist Schnitt.
  \item Gegenereignisse sparen oft Rechenarbeit.
  \item Laplace nur bei Gleichverteilung.
  \item Bei gegebener Wahrscheinlichkeitsfunktion immer ueber Summen arbeiten.
  \item Wahrscheinlichkeit $0$ bedeutet nicht automatisch ``unmoeglich'' in jedem allgemeinen Modell.
  \item Wahrscheinlichkeit $1$ bedeutet nicht automatisch ``sicher'' im logischen Sinn.
\end{itemize}

\section{Weitere Kapitel}
Hier kommen die naechsten Uebungsblaetter rein. Ich wuerde pro Kapitel immer diese Struktur beibehalten:
\begin{itemize}
  \item Grundbegriffe
  \item wichtigste Formeln
  \item Standard-Notation
  \item Aufgabentypen und Merksaetze
  \item typische Fallen
\end{itemize}

\section{Kapitel 1.2: Bernoulli- und Laplace-Experimente; Kombinatorik}

\subsection{Bernoulli-Experiment}
Ein Bernoulli-Experiment ist ein Zufallsexperiment mit genau zwei moeglichen Ergebnissen.
Typisch:
\[
\Omega = \{0,1\}
\]

Meist bezeichnet man:
\[
P(1) = p, \qquad P(0) = q = 1-p
\]
oder umgekehrt. Das muss in der Aufgabe klar festgelegt werden.

\textbf{Typische Deutung}
\begin{itemize}
  \item $1$ = Treffer, Erfolg, defekt, Kopf, ja
  \item $0$ = Niete, Misserfolg, in Ordnung, Zahl, nein
\end{itemize}

\textbf{Merke:}
Bei Bernoulli immer zuerst klar sagen, was als ``Erfolg'' codiert wird. Sonst baust du spaeter leicht $p$ und $q$ falsch ein.

\subsection{Laplace-Experiment}
Ein Laplace-Experiment liegt vor, wenn:
\begin{itemize}
  \item $\Omega$ endlich ist und
  \item alle Elementarergebnisse gleich wahrscheinlich sind.
\end{itemize}

Dann gilt:
\[
P(A) = \frac{|A|}{|\Omega|}
\]

\textbf{Typische Beispiele}
\begin{itemize}
  \item faire Muenze
  \item fairer Wuerfel
  \item Lottoziehungen
  \item Pokerblaetter
  \item zufaellige Reihenfolgen / Permutationen
\end{itemize}

\textbf{Merke:}
Bei Laplace-Aufgaben brauchst du fast immer nur:
\begin{enumerate}
  \item Anzahl aller moeglichen Ergebnisse
  \item Anzahl der guenstigen Ergebnisse
\end{enumerate}
Dann einfach Quotient bilden.

\subsection{Kombinatorik: die wichtigste Schubladenregel}
Bevor du eine Formel benutzt, frage immer:
\begin{enumerate}
  \item Wird eine \textbf{Reihenfolge} beachtet?
  \item Sind \textbf{Wiederholungen} erlaubt?
\end{enumerate}

Dann landet die Aufgabe in genau einer von vier Schubladen:

\begin{center}
\begin{tabular}{|c|c|c|}
\hline
 & ohne Wiederholung & mit Wiederholung \\
\hline
Reihenfolge wichtig & Variation & Variation \\
\hline
Reihenfolge egal & Kombination & Kombination mit Wiederholung \\
\hline
\end{tabular}
\end{center}

\textbf{Zusaetzlich:} Wenn alle $n$ Elemente angeordnet werden, spricht man von einer Permutation.

\subsection{Permutationen}
\textbf{1. Alle Elemente verschieden}

Anzahl aller Anordnungen von $n$ verschiedenen Elementen:
\[
\operatorname{Perm}(n) = n!
\]

\textbf{2. Elemente in Klassen, innerhalb der Klasse nicht unterscheidbar}

Sind $n$ Elemente in Klassen mit Groessen $k_1, k_2, \dots, k_r$ eingeteilt, dann:
\[
\operatorname{Perm}(n;k_1,\dots,k_r) = \frac{n!}{k_1!\cdot k_2!\cdots k_r!}
\]
mit
\[
k_1 + k_2 + \dots + k_r = n
\]

\textbf{Typische Beispiele}
\begin{itemize}
  \item Wortanordnungen wie \texttt{MISSISSIPPI}
  \item Sitzordnungen nach Klassen wie Mann/Frau
\end{itemize}

\textbf{Merke:}
Wenn gleiche Objekte mehrfach vorkommen, erst $n!$ denken und dann durch die inneren Vertauschungen teilen.

\subsection{Variationen: Auswahl mit Reihenfolge}
\textbf{1. Ohne Wiederholung}

Anzahl geordneter $k$-Tupel aus $n$ Elementen:
\[
V(n,k) = n\cdot(n-1)\cdot \dots \cdot (n-k+1) = \frac{n!}{(n-k)!}
\qquad (k \leq n)
\]

\textbf{Deutung:}
Ziehen ohne Zuruecklegen, Reihenfolge wichtig.

\textbf{2. Mit Wiederholung}

Anzahl geordneter $k$-Tupel aus $n$ Elementen:
\[
V(n,k) = n^k
\]

\textbf{Deutung:}
Ziehen mit Zuruecklegen, Reihenfolge wichtig.

\textbf{Typische Beispiele}
\begin{itemize}
  \item PIN-Codes
  \item Folgen von Wuerfelergebnissen
  \item Folgen von Geburtstagen
\end{itemize}

\subsection{Kombinationen: Auswahl ohne Reihenfolge}
\textbf{1. Ohne Wiederholung}

Anzahl $k$-elementiger Teilmengen aus $n$ Elementen:
\[
C(n,k) = \binom{n}{k} = \frac{n!}{k!(n-k)!}
\qquad (k \leq n)
\]

\textbf{Deutung:}
Ziehen ohne Zuruecklegen, Reihenfolge egal.

\textbf{Typische Beispiele}
\begin{itemize}
  \item Lotto ``6 aus 49''
  \item Pokerblatt mit 5 Karten
  \item Spiele ``jede Mannschaft gegen jede andere''
\end{itemize}

\textbf{2. Mit Wiederholung}

Anzahl $k$-elementiger Multimengen aus $n$ Elementen:
\[
C(n,k) = \binom{n-1+k}{k}
\]

\textbf{Deutung:}
Ziehen mit Zuruecklegen, Reihenfolge egal.

\textbf{Typische Beispiele}
\begin{itemize}
  \item Stimmverteilungen
  \item Auswahl mit Mehrfachnennungen
\end{itemize}

\subsection{Potenzmenge}
Fuer jede endliche Menge $M$ gilt:
\[
|\mathcal{P}(M)| = 2^{|M|}
\]

\textbf{Merke:}
Jedes Element ist entweder drin oder nicht drin. Also pro Element genau $2$ Moeglichkeiten.

\subsection{Wichtige Standardformeln auf einen Blick}
\begin{align*}
\operatorname{Perm}(n) &= n! \\
\operatorname{Perm}(n;k_1,\dots,k_r) &= \frac{n!}{k_1!\cdots k_r!} \\
V(n,k) \text{ ohne Wdh.} &= \frac{n!}{(n-k)!} \\
V(n,k) \text{ mit Wdh.} &= n^k \\
C(n,k) \text{ ohne Wdh.} &= \binom{n}{k} = \frac{n!}{k!(n-k)!} \\
C(n,k) \text{ mit Wdh.} &= \binom{n-1+k}{k}
\end{align*}

\subsection{Wann nehme ich was?}
\begin{itemize}
  \item \textbf{Permutation:} alle Objekte werden angeordnet.
  \item \textbf{Variation:} nur $k$ aus $n$, Reihenfolge wichtig.
  \item \textbf{Kombination:} nur $k$ aus $n$, Reihenfolge egal.
  \item \textbf{Mit Wiederholung:} gleiches Element darf mehrfach vorkommen.
  \item \textbf{Ohne Wiederholung:} ein gezogenes Element ist danach weg.
\end{itemize}

\textbf{Klausurregel:}
Wenn du unsicher bist, schreib dir ein Mini-Beispiel mit $n=3$, $k=2$ hin und ueberlege, ob z.B. $(a,b)$ und $(b,a)$ verschieden sein sollen.

\subsection{Laplace-Aufgaben mit Kombinatorik}
Oft gilt:
\[
P(A) = \frac{|A|}{|\Omega|}
\]
aber diesmal werden $|A|$ und $|\Omega|$ mit Kombinatorik bestimmt.

\textbf{Lotto ``6 aus 49''}
\[
|\Omega| = \binom{49}{6}
\]

Fuer ``genau $4$ Richtige'':
\[
|A| = \binom{6}{4}\binom{43}{2}
\]
also
\[
P(\text{genau 4 Richtige}) = \frac{\binom{6}{4}\binom{43}{2}}{\binom{49}{6}}
\]

\textbf{Merke:}
``genau $k$ Richtige'' bedeutet fast immer:
\[
\frac{\binom{\text{eigene Trefferzahlen}}{k}\binom{\text{restliche falsche}}{\text{Rest}}}{\binom{\text{gesamt}}{\text{Ziehungszahl}}}
\]

\subsection{Geburtstagsparadoxon}
Gesucht:
\[
A = \{\text{mindestens zwei Personen haben am gleichen Tag Geburtstag}\}
\]

Leichter ist das Gegenereignis:
\[
\overline{A} = \{\text{alle Geburtstage verschieden}\}
\]

Fuer $k \leq 365$ gilt:
\[
P(A) = 1 - P(\overline{A})
      = 1 - \frac{V(365,k)}{365^k}
      = 1 - \frac{365\cdot 364 \cdot \dots \cdot (365-k+1)}{365^k}
\]

\textbf{Merke:}
Bei ``mindestens zwei gleich'' fast immer ueber das Gegenereignis ``alle verschieden'' gehen.

\subsection{Ein- und Ausschlussprinzip}
Fuer endliche Mengen:
\[
|A \cup B| = |A| + |B| - |A \cap B|
\]

Fuer drei Mengen:
\[
|A \cup B \cup C|
= |A| + |B| + |C|
- |A \cap B| - |A \cap C| - |B \cap C|
+ |A \cap B \cap C|
\]

Allgemein: abwechselnd addieren und subtrahieren.

\textbf{Analog fuer Wahrscheinlichkeiten}
\[
P(A \cup B) = P(A) + P(B) - P(A \cap B)
\]

\[
P(A \cup B \cup C)
= P(A) + P(B) + P(C)
- P(A \cap B) - P(A \cap C) - P(B \cap C)
+ P(A \cap B \cap C)
\]

\textbf{Merke:}
Wenn nach ``mindestens eines der Ereignisse tritt ein'' gefragt wird und die Ereignisse sich ueberlappen koennen, ist das Ein- und Ausschlussprinzip oft die richtige Waffe.

\subsection{Typische Aufgabentypen aus dem Blatt}
\textbf{Multiple Choice mit Raten}

Pro Frage:
\[
p = \frac{1}{3}, \qquad q = \frac{2}{3}
\]

Fuer genau $k$ richtige bei $10$ Fragen:
\[
\binom{10}{k}\left(\frac{1}{3}\right)^k\left(\frac{2}{3}\right)^{10-k}
\]

Fuer ``mindestens 8 richtig'':
\[
\sum_{k=8}^{10} \binom{10}{k}\left(\frac{1}{3}\right)^k\left(\frac{2}{3}\right)^{10-k}
\]

\textbf{Personen in Reihenfolge}

Wenn alle Reihenfolgen gleich wahrscheinlich sind:
\[
|\Omega| = n!
\]

Bei guenstigen Faellen oft mit ``Blockbildung'' arbeiten:
z.B. Ehepaar nebeneinander $\Rightarrow$ Paar als ein Block zaehlen.

\textbf{Mehrfach wuerfeln und alle Ergebnisse verschieden}

Bei $k$ Wuerfen mit fairem Wuerfel:
\[
|\Omega| = 6^k
\]

Falls alle $k$ Ergebnisse verschieden sein sollen:
\[
|A| = V(6,k) = \frac{6!}{(6-k)!}
\]

also
\[
P(A) = \frac{V(6,k)}{6^k}
\]

\subsection{Notation fuer saubere Loesungen}
Gute Standardsaetze:
\begin{itemize}
  \item ``Da die Reihenfolge relevant ist, betrachten wir Variationen.''
  \item ``Da ohne Zuruecklegen gezogen wird, sind keine Wiederholungen moeglich.''
  \item ``Da nur die Auswahl, nicht aber die Reihenfolge zaehlt, verwenden wir Kombinationen.''
  \item ``Da es sich um ein Laplace-Experiment handelt, gilt $P(A)=|A|/|\Omega|$. ''
  \item ``Wir berechnen zunaechst die Anzahl aller moeglichen und dann die Anzahl der guenstigen Ergebnisse.''
  \item ``Es ist einfacher, das Gegenereignis zu betrachten.''
\end{itemize}

\subsection{Typische Fallen}
\begin{itemize}
  \item Reihenfolge versehentlich mitgezaehlt, obwohl sie egal ist.
  \item Reihenfolge vergessen, obwohl sie wichtig ist.
  \item Mit/ohne Wiederholung verwechselt.
  \item Bei Laplace direkt Formeln hingeschrieben, ohne $|\Omega|$ und $|A|$ sauber zu definieren.
  \item ``mindestens eins'' nicht ueber Gegenereignis oder Ein-/Ausschlussprinzip gerechnet.
  \item Lotto/Poker als Tupel statt als Mengen modelliert.
\end{itemize}

\section{Kapitel 1.3: Zufallsvariablen und zufaellige Elemente}

\subsection{Grundidee}
Oft interessiert beim Zufallsexperiment nicht das rohe Ergebnis $\omega \in \Omega$, sondern nur ein bestimmtes Merkmal davon.

\textbf{Beispiele}
\begin{itemize}
  \item Bei zwei Wuerfeln interessiert nur die Summe.
  \item Bei zwei Wuerfeln interessiert nur das Produkt.
  \item Bei zwei Kindern interessiert nur die Geschlechterzusammensetzung.
  \item Bei drei nicht unterscheidbaren Wuerfeln interessiert nur die Multimenge der Augenzahlen.
\end{itemize}

\textbf{Merke:}
Man startet mit einem sauberen Wahrscheinlichkeitsraum $(\Omega,P)$ und bildet daraus ueber eine Abbildung $X$ ein neues Modell fuer das interessierende Merkmal.

\subsection{Zufallsvariable}
Eine Zufallsvariable ist eine Abbildung
\[
X : \Omega \to \mathbb{R}
\]

\textbf{Bedeutung:}
\begin{itemize}
  \item jedem Ergebnis $\omega$ wird ein Zahlenwert $X(\omega)$ zugeordnet,
  \item dieser Zahlenwert ist das Merkmal, das uns interessiert.
\end{itemize}

\textbf{Typische Beispiele}
\begin{itemize}
  \item Augensumme zweier Wuerfel
  \item Augenprodukt zweier Wuerfel
  \item Koerpergroesse
  \item Studiensemester
\end{itemize}

\subsection{\texorpdfstring{Verteilung von $X$ / Bildmass}{Verteilung von X / Bildmass}}
Aus $X$ entsteht eine neue Verteilung auf der Wertemenge von $X$:
\[
P_X(A) := P(X \in A)
\]

gleichbedeutend mit
\[
P_X(A) = P(X^{-1}(A))
\]

Dabei ist
\[
X^{-1}(A) = \{\omega \in \Omega \mid X(\omega) \in A\}
\]
das Urbild von $A$.

\textbf{Wichtige Schreibweisen}
\begin{align*}
\{X \in A\} &= X^{-1}(A) \\
\{X = a\} &= \{\omega \in \Omega \mid X(\omega)=a\} \\
\{X \le a\} &= \{\omega \in \Omega \mid X(\omega)\le a\} \\
\{X \ge a\} &= \{\omega \in \Omega \mid X(\omega)\ge a\}
\end{align*}

\textbf{Merke:}
Wenn du eine Wahrscheinlichkeit fuer die Zufallsvariable suchst, rechnest du immer im urspruenglichen Raum $\Omega$ ueber das passende Urbild.

\subsection{Standardrezept fuer Aufgaben}
\begin{enumerate}
  \item Urspruenglichen Wahrscheinlichkeitsraum $(\Omega,P)$ angeben.
  \item Zufallsvariable $X$ definieren.
  \item Gesuchtes Ereignis als $\{X \in A\}$ oder $\{X=a\}$ schreiben.
  \item Urbild in $\Omega$ bestimmen.
  \item Wahrscheinlichkeit dort ausrechnen.
\end{enumerate}

\subsection{Beispiel: Augensumme zweier fairer Wuerfel}
Ursprungsraum:
\[
\Omega = \{1,2,3,4,5,6\}^2
\]
mit Gleichverteilung.

Zufallsvariable:
\[
X(i,j)=i+j
\]

Wertemenge:
\[
\Omega_X=\{2,3,\dots,12\}
\]

Beispiel fuer $A=\{2,3\}$:
\[
X^{-1}(A)=\{(1,1),(1,2),(2,1)\}
\]

also
\[
P_X(A)=P(X\in A)=\frac{3}{36}=\frac{1}{12}
\]

\textbf{Merke:}
Die Summe ist nicht gleichverteilt auf $\{2,\dots,12\}$, weil verschiedene Summen unterschiedlich viele Urbilder haben.

\subsection{Beispiel: Augenprodukt zweier fairer Wuerfel}
Zufallsvariable:
\[
X(i,j)=i\cdot j
\]

Gesucht z.B.:
\[
A=\{\text{Produkt ist durch } 6 \text{ teilbar}\}
\]

Dann:
\[
P_X(A)=P(X\in A)
\]

Vorgehen:
\begin{itemize}
  \item alle Paare $(i,j)$ suchen, fuer die $i\cdot j$ durch $6$ teilbar ist,
  \item diese Paare zaehlen,
  \item durch $36$ teilen.
\end{itemize}

Im Blatt:
\[
P(X\in A)=\frac{15}{36}=\frac{5}{12}
\]

\subsection{Zufaelliges Element}
Allgemeiner als eine Zufallsvariable:
\[
X : \Omega \to \Omega'
\]
wobei $\Omega'$ irgendeine Menge sein darf.

Dann heisst $X$ ein zufaelliges Element.

\textbf{Unterschied}
\begin{itemize}
  \item Zufallsvariable: Zielmenge ist numerisch.
  \item Zufaelliges Element: Zielmenge kann auch aus Mengen, Multimengen, Kartenblaettern, Worten usw. bestehen.
\end{itemize}

\textbf{Beziehung}
\begin{itemize}
  \item Jede Zufallsvariable ist ein zufaelliges Element.
  \item Nicht jedes zufaellige Element ist eine Zufallsvariable.
\end{itemize}

\subsection{Beispiel: drei nicht unterscheidbare Wuerfel}
Wenn die Wuerfel nicht unterscheidbar sind, ist ein Tupel wie $(3,2,6)$ kein gutes Endmodell, weil die Reihenfolge bedeutungslos ist.

Besser:
\begin{itemize}
  \item zuerst sauberes Modell mit unterscheidbaren Wuerfeln:
  \[
  \Omega=\{1,2,3,4,5,6\}^3
  \]
  mit Gleichverteilung,
  \item dann Abbildung auf die passende Struktur:
  \[
  X(i,j,k)=\text{Multimenge } \{i,j,k\}
  \]
\end{itemize}

So bekommt man die richtige Verteilung auf den Multimengen.

Beispiel:
\[
X(i,j,k)=\{2,2,3\}
\]
ist genau fuer
\[
(2,2,3), (2,3,2), (3,2,2)
\]
erfuellt. Also:
\[
P(X=\{2,2,3\})=\frac{3}{6^3}=\frac{3}{216}=\frac{1}{72}
\]

\textbf{Merke:}
Wenn du von einem feineren Modell auf ein groberes Modell abbildest, ist die neue Verteilung meistens nicht gleichverteilt.

\subsection{Extrem wichtige Klausurfalle}
Nur weil die neue Ergebnismenge klein und symmetrisch aussieht, darfst du nicht automatisch Gleichverteilung annehmen.

\textbf{Beispiel Kinder-Geschlechter}

Feines Modell:
\[
\Omega=\{(0,0),(1,0),(0,1),(1,1)\}
\]
mit Gleichverteilung, etwa
\[
1=\text{Maedchen}, \qquad 0=\text{Junge}
\]

Ereignis ``mindestens ein Junge'':
\[
S=\{(0,0),(1,0),(0,1)\}
\]
also
\[
P(S)=\frac{3}{4}
\]

Grobes Modell ueber Multimengen:
\[
\{0,0\}, \{0,1\}, \{1,1\}
\]

Diese drei Ergebnisse sind nicht gleich wahrscheinlich, denn
\begin{align*}
P(\{0,0\}) &= \frac{1}{4} \\
P(\{0,1\}) &= \frac{2}{4}=\frac{1}{2} \\
P(\{1,1\}) &= \frac{1}{4}
\end{align*}

Deshalb ist
\[
P(S)=P(\{0,0\})+P(\{0,1\})=\frac{1}{4}+\frac{1}{2}=\frac{3}{4}
\]
und nicht $\frac{2}{3}$.

\textbf{Merke:}
Wenn mehrere Ergebnisse des feinen Modells auf dasselbe grobe Ergebnis fallen, dann bekommt dieses grobe Ergebnis mehr Gewicht.

\subsection{Chevalier-de-Mere-Typ-Aufgaben}
Bei drei Wuerfeln und Summen wie $11$ oder $12$ reicht es nicht, nur Zerlegungen als Summen zu vergleichen.

Warum?
\begin{itemize}
  \item Dieselbe Zerlegung kann unterschiedlich viele geordnete Tupel erzeugen.
  \item Beispiel: $5+5+1$ liefert weniger verschiedene Tupel als $6+4+1$.
\end{itemize}

\textbf{Klausurregel:}
Wenn das Grundmodell geordnete Tupel sind, musst du auch die Anzahl der geordneten Tupel zaehlen.

\subsection{Wann benutze ich Zufallsvariable vs. direktes Ereignis?}
\begin{itemize}
  \item Direktes Ereignis: wenn im Originalraum sowieso leicht gezaehlt werden kann.
  \item Zufallsvariable: wenn sich die Frage natuerlich ueber ein Merkmal wie Summe, Produkt, Anzahl, Gewicht, Groesse formulieren laesst.
\end{itemize}

\textbf{Praktisch:}
In Rechnungen wird trotzdem meist wieder im Originalraum ueber das Urbild gezaehlt.

\subsection{Notation fuer saubere Loesungen}
Gute Standardsaetze:
\begin{itemize}
  \item ``Wir definieren die Zufallsvariable $X:\Omega \to \mathbb{R}$ durch $X(\omega)=\dots$. ''
  \item ``Gesucht ist $P(X \in A)$. ''
  \item ``Dazu bestimmen wir zunaechst das Urbild $X^{-1}(A)$. ''
  \item ``Da auf $\Omega$ Gleichverteilung herrscht, gilt $P(X \in A)=|X^{-1}(A)|/|\Omega|$. ''
  \item ``Die grobe Ergebnismenge ist nicht gleichverteilt; die Verteilung muss aus dem Ursprungsraum transportiert werden.''
\end{itemize}

\subsection{Mini-Merkliste}
\begin{itemize}
  \item Zufallsvariable = numerisches Merkmal eines Zufallsexperiments.
  \item Zufaelliges Element = allgemeiner; Zielmenge muss nicht numerisch sein.
  \item Wahrscheinlichkeiten von $X$ immer ueber Urbilder im urspruenglichen Raum berechnen.
  \item Bildverteilungen sind oft nicht gleichverteilt, auch wenn der Ursprungsraum gleichverteilt ist.
  \item Bei vergroebberten Modellen entscheidet die Anzahl der Urbilder ueber die Wahrscheinlichkeiten.
  \item Zerlegungen zaehlen reicht nicht, wenn das Grundmodell geordnete Tupel hat.
\end{itemize}

\section{Kapitel 1.4: Bedingte Wahrscheinlichkeiten}

\subsection{Grundidee}
Eine bedingte Wahrscheinlichkeit beschreibt, wie sich Wahrscheinlichkeiten aendern, wenn man zusaetzliche Information kennt.

\textbf{Beispiel}
\begin{itemize}
  \item $A$ = ``Augenzahl ist gerade''
  \item $B$ = ``Augenzahl ist mindestens 4''
\end{itemize}

Ohne Zusatzwissen gilt:
\[
P(A)=\frac{3}{6}=\frac{1}{2}
\]

Wenn aber bekannt ist, dass $B$ eingetreten ist, dann bleibt nur
\[
B=\{4,5,6\}
\]
und davon sind fuer $A$ guenstig:
\[
A \cap B=\{4,6\}
\]

Also:
\[
P(A\mid B)=\frac{2}{3}
\]

\subsection{Definition}
Fuer Ereignisse $A,B$ mit $P(B)>0$ gilt:
\[
P(A\mid B)=\frac{P(A \cap B)}{P(B)}
\]

\textbf{Bedeutung:}
``Wahrscheinlichkeit von $A$ unter der Bedingung, dass $B$ bereits eingetreten ist.''

\textbf{Merke:}
Im Nenner steht immer die Bedingung.

\subsection{Wichtige Umformung}
Aus der Definition folgt direkt:
\[
P(A \cap B)=P(A\mid B)\cdot P(B)
\]

und genauso
\[
P(A \cap B)=P(B\mid A)\cdot P(A)
\]

Das ist der \textbf{Multiplikationssatz}.

\subsection{Multiplikationssatz}
Fuer zwei Ereignisse:
\[
P(A \cap B)=P(A\mid B)\cdot P(B)=P(B\mid A)\cdot P(A)
\]

Fuer drei Ereignisse:
\[
P(A_1 \cap A_2 \cap A_3)
=P(A_3\mid A_1 \cap A_2)\cdot P(A_2\mid A_1)\cdot P(A_1)
\]

\textbf{Merke:}
Bei Pfaden im Baum immer entlang des Pfades multiplizieren.

\subsection{Ereignisbaeume}
Ein Ereignisbaum ist hilfreich bei mehrstufigen Zufallsexperimenten.

\textbf{Regeln}
\begin{itemize}
  \item An jeder Kante steht eine bedingte Wahrscheinlichkeit.
  \item Die Wahrscheinlichkeit eines vollstaendigen Pfades bekommt man durch Multiplikation.
  \item Wenn mehrere disjunkte Pfade zum gesuchten Ereignis fuehren, addiert man deren Pfadwahrscheinlichkeiten.
\end{itemize}

\textbf{Klausurregel}
\begin{itemize}
  \item \textbf{entlang eines Pfades:} multiplizieren
  \item \textbf{zwischen verschiedenen guenstigen Pfaden:} addieren
\end{itemize}

\subsection{Satz von der totalen Wahrscheinlichkeit}
Ist
\[
\Omega = B_1 \cup B_2 \cup \dots \cup B_n
\]
eine Zerlegung in paarweise disjunkte Ereignisse mit $P(B_i)>0$, dann gilt fuer jedes Ereignis $A$:
\[
P(A)=\sum_{i=1}^n P(A\mid B_i)\,P(B_i)
\]

auch geschrieben als
\[
P(A)=P(A\cap B_1)+\dots+P(A\cap B_n)
\]

\textbf{Wann benutzen?}
\begin{itemize}
  \item Wenn es mehrere moegliche Ursachen/Faelle $B_i$ gibt.
  \item Wenn $P(A)$ insgesamt gesucht ist, aber nur die bedingten Wahrscheinlichkeiten $P(A\mid B_i)$ bekannt sind.
\end{itemize}

\textbf{Merke:}
Totale Wahrscheinlichkeit = \emph{vorwaerts} rechnen.

\subsection{Satz von Bayes}
Unter denselben Voraussetzungen gilt fuer $P(A)>0$:
\[
P(B_i\mid A)=\frac{P(A\mid B_i)\cdot P(B_i)}{P(A)}
\]

Mit totaler Wahrscheinlichkeit im Nenner:
\[
P(B_i\mid A)=
\frac{P(A\mid B_i)\cdot P(B_i)}
{\sum_{j=1}^n P(A\mid B_j)\cdot P(B_j)}
\]

\textbf{Wann benutzen?}
\begin{itemize}
  \item Wenn ein Ergebnis $A$ beobachtet wurde
  \item und du danach wissen willst, wie wahrscheinlich eine bestimmte Ursache $B_i$ war
\end{itemize}

\textbf{Merke:}
Bayes = \emph{rueckwaerts} rechnen.

\subsection{Unabhaengigkeit}
Zwei Ereignisse $A$ und $B$ heissen unabhaengig, wenn
\[
P(A \cap B)=P(A)\cdot P(B)
\]

Falls $P(B)>0$, ist das aequivalent zu
\[
P(A\mid B)=P(A)
\]

\textbf{Bedeutung:}
Das Wissen, dass $B$ eingetreten ist, aendert die Wahrscheinlichkeit von $A$ nicht.

\textbf{Merke:}
Unabhaengig heisst \emph{nicht} disjunkt.

\textbf{Typischer Gegencheck}
\begin{itemize}
  \item disjunkt mit positiven Wahrscheinlichkeiten $\Rightarrow$ nicht unabhaengig
  \item unabhaengig $\Rightarrow$ Schnitt ist normalerweise nicht leer
\end{itemize}

\subsection{Vierfeldertafel}
Bei Ja/Nein-Situationen mit zwei Merkmalen ist eine Vierfeldertafel oft am saubersten:

\[
\begin{array}{c|cc|c}
 & T_0 & T_1 & \text{Summe} \\
\hline
K_0 & P(K_0\cap T_0) & P(K_0\cap T_1) & P(K_0) \\
K_1 & P(K_1\cap T_0) & P(K_1\cap T_1) & P(K_1) \\
\hline
\text{Summe} & P(T_0) & P(T_1) & 1
\end{array}
\]

\textbf{Merke:}
Zeilen- und Spaltensummen muessen immer passen. Das ist auch ein guter Fehlercheck.

\subsection{Typische Standardaufgaben}
\textbf{1. Kanal / Uebertragungsfehler}

Wenn z.B.
\begin{itemize}
  \item $B_0$ = ``0 wurde gesendet''
  \item $B_1$ = ``1 wurde gesendet''
  \item $A$ = ``0 wurde empfangen''
\end{itemize}

dann:
\[
P(A)=P(A\mid B_0)P(B_0)+P(A\mid B_1)P(B_1)
\]

\textbf{2. Hersteller / Defekt}

Wenn $B_i$ = ``Produkt stammt von Hersteller $i$'' und $D$ = ``Produkt ist defekt'', dann:
\[
P(D)=\sum_i P(D\mid B_i)P(B_i)
\]

und danach
\[
P(B_i\mid D)=\frac{P(D\mid B_i)P(B_i)}{P(D)}
\]

\textbf{3. Treffer und unbekannte Ursache}

Wenn mehrere Schuetzen unterschiedlich haeufig schiessen und unterschiedlich gut treffen:
\begin{itemize}
  \item erst Priors $P(B_i)$ aus den Schussanteilen bestimmen
  \item dann Trefferwahrscheinlichkeit $P(T\mid B_i)$ einsetzen
  \item danach Bayes fuer $P(B_i\mid T)$
\end{itemize}

\textbf{4. Medizinischer Test}

Typische Notation:
\begin{align*}
K_1 &= \text{krank} \\
K_0 &= \text{gesund} \\
T_1 &= \text{Test positiv} \\
T_0 &= \text{Test negativ}
\end{align*}

Dann:
\begin{align*}
P(T_1\mid K_1) &= \text{Sensitivitaet} \\
P(T_0\mid K_0) &= \text{Spezifitaet}
\end{align*}

Gesucht ist oft:
\[
P(K_1\mid T_1)
=\frac{P(T_1\mid K_1)P(K_1)}
{P(T_1\mid K_1)P(K_1)+P(T_1\mid K_0)P(K_0)}
\]

\subsection{Testparadoxon}
Der Denkfehler:
\begin{itemize}
  \item viele verwechseln $P(T_1\mid K_1)$ mit $P(K_1\mid T_1)$
  \item das ist fast nie dasselbe
\end{itemize}

\textbf{Merke brutal wichtig:}
\[
P(\text{positiv}\mid \text{krank}) \neq P(\text{krank}\mid \text{positiv})
\]

Der zweite Wert haengt stark von der Grundhaeufigkeit der Krankheit ab:
\[
P(K_1)
\]

Wenn die Krankheit selten ist, koennen trotz gutem Test viele positive Ergebnisse falsch positiv sein.

\subsection{Saubere Standardnotation}
Gute Standardsaetze:
\begin{itemize}
  \item ``Wir betrachten die Zerlegung $\Omega=B_1\cup\dots\cup B_n$. ''
  \item ``Mit dem Satz von der totalen Wahrscheinlichkeit gilt \dots''
  \item ``Mit dem Satz von Bayes folgt \dots''
  \item ``Da $P(B)>0$, ist die bedingte Wahrscheinlichkeit $P(A\mid B)$ definiert.''
  \item ``Die Pfadwahrscheinlichkeit ergibt sich als Produkt der Kantenwahrscheinlichkeiten.''
  \item ``Das gesuchte Ereignis setzt sich aus paarweise disjunkten Teilfaellen zusammen; daher addieren wir die Wahrscheinlichkeiten.''
\end{itemize}

\subsection{Typische Fallen}
\begin{itemize}
  \item $P(A\mid B)$ mit $P(B\mid A)$ verwechselt.
  \item Im Nenner der bedingten Wahrscheinlichkeit das falsche Ereignis genommen.
  \item Im Baum addiert, obwohl multipliziert werden muss, oder umgekehrt.
  \item Bayes benutzt, obwohl eigentlich nur totale Wahrscheinlichkeit gebraucht wird.
  \item Grundwahrscheinlichkeit $P(B_i)$ vergessen.
  \item Bei Testaufgaben die Grundhaeufigkeit der Krankheit ignoriert.
  \item Unabhaengigkeit und Disjunktheit verwechselt.
\end{itemize}

\subsection{Mini-Merkliste}
\begin{itemize}
  \item Bedingte Wahrscheinlichkeit: $\displaystyle P(A\mid B)=\frac{P(A\cap B)}{P(B)}$.
  \item Multiplikationssatz: Schnitt = bedingte Wkt. mal Bedingung.
  \item Totale Wahrscheinlichkeit = mehrere Ursachen vorwaerts aufsummieren.
  \item Bayes = beobachtetes Ergebnis rueckwaerts auf Ursachen verteilen.
  \item Im Baum: entlang multiplizieren, ueber disjunkte Pfade addieren.
  \item Testparadoxon: $P(T_1\mid K_1)$ ist nicht das, was meist gefragt ist.
\end{itemize}

\section{Kapitel 1.5: Unabhaengigkeit von Ereignissen und Zufallsvariablen}

\subsection{Unabhaengigkeit von Ereignissen}
Zwei Ereignisse $A,B$ heissen stochastisch unabhaengig genau dann, wenn
\[
P(A \cap B)=P(A)\cdot P(B)
\]

\textbf{Bedeutung:}
Das Eintreten von $B$ aendert die Wahrscheinlichkeit von $A$ nicht, und umgekehrt.

Falls $P(B)>0$, ist das aequivalent zu
\[
P(A\mid B)=P(A)
\]

Falls $P(A)>0$, ist das aequivalent zu
\[
P(B\mid A)=P(B)
\]

\textbf{Merke:}
Am sichersten prueft man Unabhaengigkeit fast immer mit
\[
P(A \cap B)\stackrel{?}{=}P(A)P(B)
\]
weil das auch dann funktioniert, wenn eine Wahrscheinlichkeit $0$ ist.

\subsection{Wichtige aequivalente Aussagen}
Unter geeigneten Voraussetzungen sind aequivalent:
\begin{align*}
P(A \cap B)&=P(A)P(B) \\
P(A\mid B)&=P(A) \\
P(B\mid A)&=P(B) \\
P(A\mid \overline{B})&=P(A) \\
P(\overline{A}\mid B)&=P(\overline{A})
\end{align*}

\textbf{Merke:}
Unabhaengigkeit bleibt auch bei Gegenereignissen erhalten:
\begin{itemize}
  \item aus $A \perp B$ folgt auch $A \perp \overline{B}$
  \item aus $A \perp B$ folgt auch $\overline{A} \perp B$
  \item aus $A \perp B$ folgt auch $\overline{A} \perp \overline{B}$
\end{itemize}

\subsection{Extrem wichtige Abgrenzung}
\textbf{Unabhaengig} ist nicht dasselbe wie \textbf{disjunkt}.

\textbf{Disjunkt:}
\[
A \cap B = \emptyset
\]

\textbf{Unabhaengig:}
\[
P(A \cap B)=P(A)P(B)
\]

Falls $A$ und $B$ disjunkt sind und beide positive Wahrscheinlichkeit haben, dann gilt
\[
P(A \cap B)=0 \neq P(A)P(B)
\]
also sind sie \textbf{nicht} unabhaengig.

\textbf{Merke:}
Disjunkt mit positiver Wahrscheinlichkeit bedeutet fast immer abhaengig.

\subsection{Spezialfaelle}
\begin{itemize}
  \item Wenn $P(A)=0$ oder $P(B)=0$, dann sind $A$ und $B$ unabhaengig.
  \item Wenn $P(A)=1$ oder $P(B)=1$, dann sind $A$ und $B$ ebenfalls unabhaengig.
\end{itemize}

\textbf{Warum?}
Dann passt die Gleichung
\[
P(A \cap B)=P(A)P(B)
\]
automatisch.

\subsection{Ereignis unabhaengig von sich selbst}
Ein Ereignis $A$ ist genau dann unabhaengig von sich selbst, wenn
\[
P(A \cap A)=P(A)^2
\]
also
\[
P(A)=P(A)^2
\]

Daraus folgt:
\[
P(A)\in\{0,1\}
\]

\textbf{Merke:}
Ein Ereignis ist nur dann von sich selbst unabhaengig, wenn es fast unmoeglich oder sicher ist.

\subsection{Typische Pruefung von Unabhaengigkeit in Tabellen}
Bei einer Tabelle mit gemeinsamen Wahrscheinlichkeiten:
\begin{itemize}
  \item zuerst Randwahrscheinlichkeiten berechnen
  \item dann vergleichen:
  \[
  P(A \cap B)\stackrel{?}{=}P(A)P(B)
  \]
\end{itemize}

Beispielstruktur:
\[
\begin{array}{c|cc}
 & B & \overline{B} \\
\hline
A & P(A \cap B) & P(A \cap \overline{B}) \\
\overline{A} & P(\overline{A} \cap B) & P(\overline{A} \cap \overline{B})
\end{array}
\]

\subsection{Unabhaengigkeit von Zufallsvariablen}
Zwei Zufallsvariablen $X,Y$ heissen unabhaengig, wenn fuer alle Teilmengen $A,B \subseteq \mathbb{R}$ gilt:
\[
P(X \in A;\, Y \in B)=P(X \in A)\cdot P(Y \in B)
\]

Bei diskreten Zufallsvariablen reicht es, fuer alle Werte $a,b$ zu pruefen:
\[
P(X=a, Y=b)=P(X=a)\cdot P(Y=b)
\]

\textbf{Merke:}
Bei diskreten Variablen prueft man meistens punktweise mit den einzelnen Werten.

\subsection{Gemeinsame Verteilung}
Die gemeinsame Verteilung von $X$ und $Y$ ist die Verteilung des zufaelligen Elements
\[
Z=(X,Y)
\]

Also:
\[
P_{X,Y}(a,b)=P(X=a,Y=b)
\]

Das ist eine Tabelle aller gemeinsamen Wahrscheinlichkeiten.

\subsection{Randverteilungen}
Aus der gemeinsamen Verteilung bekommt man die Randverteilungen durch Aufsummieren ueber Zeilen bzw. Spalten:
\begin{align*}
P(X=a) &= \sum_b P(X=a,Y=b) \\
P(Y=b) &= \sum_a P(X=a,Y=b)
\end{align*}

\textbf{Merke:}
``Randverteilung'' heisst so, weil die Summen in einer Tabelle am Rand landen.

\subsection{Kriterium ueber gemeinsame und Randverteilungen}
Zwei Zufallsvariablen $X,Y$ sind genau dann unabhaengig, wenn fuer alle $a,b$ gilt:
\[
P(X=a,Y=b)=P(X=a)\cdot P(Y=b)
\]

Also:
\[
\text{gemeinsame Verteilung} = \text{Produkt der Randverteilungen}
\]

\textbf{Praktisches Rezept}
\begin{enumerate}
  \item Gemeinsame Verteilung aufstellen.
  \item Randverteilungen ausrechnen.
  \item Fuer passende Werte $a,b$ pruefen, ob das Produkt passt.
\end{enumerate}

\subsection{Typische Beispiele}
\textbf{1. Zwei faire unterscheidbare Wuerfel}

Wenn
\begin{itemize}
  \item $X$ = Augenzahl des weissen Wuerfels
  \item $Y$ = Augenzahl des schwarzen Wuerfels
\end{itemize}
dann sind $X$ und $Y$ unabhaengig.

Denn:
\[
P(X=a)=\frac16,\qquad P(Y=b)=\frac16,\qquad P(X=a,Y=b)=\frac1{36}
\]
also
\[
P(X=a,Y=b)=P(X=a)P(Y=b)
\]

\textbf{2. Summe und einzelner Wuerfel}

Wenn
\begin{itemize}
  \item $X$ = Summe beider Augenzahlen
  \item $Y$ = Augenzahl des schwarzen Wuerfels
\end{itemize}
dann sind $X$ und $Y$ \textbf{abhaengig}.

Beispiel:
\[
P(X=2,Y=3)=0
\]
aber
\[
P(X=2)\cdot P(Y=3)=\frac1{36}\cdot\frac16 \neq 0
\]

\textbf{Merke:}
Sobald eine Variable Information ueber die andere verraet, sind sie meistens abhaengig.

\subsection{Unabhaengigkeit von $n$ Zufallsvariablen}
$X_1,\dots,X_n$ heissen vollstaendig unabhaengig, wenn fuer alle Teilmengen $A_1,\dots,A_n$ gilt:
\[
P(X_1\in A_1;\dots;X_n\in A_n)=P(X_1\in A_1)\cdots P(X_n\in A_n)
\]

\textbf{Wichtig:}
Paarweise unabhaengig ist nicht immer dasselbe wie vollstaendig unabhaengig.

\textbf{Merke:}
Wenn in der Aufgabe explizit mehrere Variablen vorkommen, genau lesen, ob wirklich ``vollstaendig unabhaengig'' gemeint ist.

\subsection{Typische Aufgabentypen}
\textbf{1. Sind zwei Ereignisse unabhaengig?}

Rezept:
\begin{enumerate}
  \item $P(A)$ berechnen
  \item $P(B)$ berechnen
  \item $P(A \cap B)$ berechnen
  \item vergleichen mit $P(A)P(B)$
\end{enumerate}

\textbf{2. Sind zwei Zufallsvariablen unabhaengig?}

Rezept:
\begin{enumerate}
  \item gemeinsame Verteilung bestimmen
  \item Randverteilungen bestimmen
  \item pruefen, ob fuer alle Werte das Produktkriterium gilt
\end{enumerate}

\textbf{3. Bedingte Wahrscheinlichkeit und Unabhaengigkeit}

Wenn $P(A\mid B)\neq P(A)$, dann sind $A$ und $B$ abhaengig.

\textbf{4. Statistische Tabellen}

Wenn z.B. Geschlecht und Testerfolg gegeben sind:
\begin{itemize}
  \item $M$ = Maedchen
  \item $B$ = bestanden
\end{itemize}

Dann pruefen:
\[
P(M \cap B)\stackrel{?}{=}P(M)P(B)
\]

oder alternativ
\[
P(B\mid M)\stackrel{?}{=}P(B)
\]

\subsection{Typische Fallen}
\begin{itemize}
  \item Disjunktheit mit Unabhaengigkeit verwechselt.
  \item Nur ein Beispielwert geprueft, obwohl alle noetig waeren.
  \item Bei Zufallsvariablen gemeinsame und Randverteilung verwechselt.
  \item Aus ``gleiche Randverteilungen'' faelschlich auf Unabhaengigkeit geschlossen.
  \item $P(A\mid B)=P(A)$ benutzt, obwohl $P(B)=0$ ist und der Ausdruck gar nicht definiert waere.
  \item Paarweise Unabhaengigkeit mit vollstaendiger Unabhaengigkeit verwechselt.
\end{itemize}

\subsection{Mini-Merkliste}
\begin{itemize}
  \item Ereignisse unabhaengig $\Leftrightarrow P(A \cap B)=P(A)P(B)$.
  \item Falls definiert: auch aequivalent zu $P(A\mid B)=P(A)$.
  \item Disjunkt ist nicht unabhaengig.
  \item Ereignis nur dann von sich selbst unabhaengig, wenn $P(A)=0$ oder $1$.
  \item Zufallsvariablen unabhaengig $\Leftrightarrow$ gemeinsame Verteilung = Produkt der Randverteilungen.
  \item Randverteilungen bekommt man durch Aufsummieren der gemeinsamen Verteilung.
\end{itemize}

\section{Kapitel 1.6: Haeufig vorkommende Experimente und Verteilungen}

\subsection{Die wichtigste Auswahlfrage}
Bevor du irgendeine Formel nimmst, frage:
\begin{enumerate}
  \item Gibt es eine \textbf{feste Anzahl} von Versuchen?
  \item Oder wartet man \textbf{bis zum ersten Erfolg}?
  \item Wird \textbf{mit Zuruecklegen / unabhaengig} gezogen?
  \item Oder \textbf{ohne Zuruecklegen / abhaengig}?
\end{enumerate}

Dann gilt fast immer:
\begin{itemize}
  \item feste Anzahl, unabhaengige Bernoulli-Versuche $\Rightarrow$ \textbf{Binomialverteilung}
  \item warten bis erster Erfolg $\Rightarrow$ \textbf{geometrische Verteilung}
  \item feste Anzahl, ohne Zuruecklegen $\Rightarrow$ \textbf{hypergeometrische Verteilung}
\end{itemize}

\subsection{$n$ unabhaengige Durchfuehrungen eines Experiments}
Hat das Einzelexperiment den Wahrscheinlichkeitsraum $(\Omega,P)$, dann beschreibt
\[
(\Omega^n,P')
\]
die $n$ unabhaengigen Wiederholungen.

Fuer Rechteckereignisse gilt:
\[
P'(A_1 \times \dots \times A_n)=P(A_1)\cdots P(A_n)
\]

\textbf{Merke:}
Unabhaengige Wiederholung bedeutet: Wahrscheinlichkeiten entlang der Durchfuehrungen multiplizieren.

\subsection{Bernoulli-Kette}
Ein Bernoulli-Experiment hat
\[
\Omega=\{0,1\}, \qquad P(1)=p,\qquad P(0)=q=1-p
\]

Eine \textbf{Bernoulli-Kette der Laenge $n$} ist die $n$-fache unabhaengige Wiederholung dieses Experiments.

Hat ein Tupel $(\omega_1,\dots,\omega_n)$ genau $k$ Erfolge ($1$en), dann gilt:
\[
P((\omega_1,\dots,\omega_n))=p^k q^{n-k}
\]

\textbf{Merke:}
Alle Folgen mit gleich vielen Erfolgen haben dieselbe Wahrscheinlichkeit.

\subsection{Binomialverteilung}
Wenn $X$ die Anzahl der Erfolge in $n$ unabhaengigen Bernoulli-Versuchen mit Erfolgswahrscheinlichkeit $p$ zaehlt, dann gilt:
\[
X \sim \operatorname{Bin}(n,p)
\]

mit
\[
P(X=k)=\binom{n}{k}p^k q^{n-k}
\qquad \text{fuer } k=0,1,\dots,n
\]

\textbf{Bedeutung der Faktoren}
\begin{itemize}
  \item $\binom{n}{k}$ = an wie vielen Stellen die $k$ Erfolge liegen koennen
  \item $p^k$ = Wahrscheinlichkeit der $k$ Erfolge
  \item $q^{n-k}$ = Wahrscheinlichkeit der restlichen $n-k$ Misserfolge
\end{itemize}

\subsection{Wann nehme ich die Binomialverteilung?}
\begin{itemize}
  \item feste Anzahl $n$ von Versuchen
  \item jeder Versuch hat genau zwei Ausgaenge: Erfolg / Misserfolg
  \item Erfolgswahrscheinlichkeit $p$ bleibt gleich
  \item Versuche sind unabhaengig
\end{itemize}

\textbf{Typische Formulierungen}
\begin{itemize}
  \item ``genau $k$ Treffer''
  \item ``hoechstens $k$ Ausschussteile''
  \item ``mindestens $k$ richtige Antworten''
  \item ``in $n$ Versuchen wie oft Erfolg''
\end{itemize}

\subsection{Wichtige Binomial-Summen}
\begin{align*}
P(X\le k) &= \sum_{i=0}^k \binom{n}{i}p^i q^{n-i} \\
P(X\ge k) &= \sum_{i=k}^n \binom{n}{i}p^i q^{n-i} \\
P(a\le X\le b) &= \sum_{i=a}^b \binom{n}{i}p^i q^{n-i}
\end{align*}

\textbf{Merke:}
Bei ``mindestens'' und ``hoechstens'' oft das Gegenereignis benutzen, wenn die Summe kuerzer wird.

\subsection{Typische Binomial-Beispiele}
\textbf{1. Multiple Choice mit Raten}

Pro Frage:
\[
p=\frac13,\qquad q=\frac23
\]

Bei $10$ Fragen:
\[
X \sim \operatorname{Bin}(10,\tfrac13)
\]

Bestanden ab mindestens $8$ richtig:
\[
P(X\ge 8)=\sum_{k=8}^{10}\binom{10}{k}\left(\frac13\right)^k\left(\frac23\right)^{10-k}
\]

\textbf{2. Ausschuss in Lieferung}

Wenn jedes Geraet mit Wahrscheinlichkeit $p=0.02$ Ausschuss ist und $100$ Geraete unabhaengig produziert werden:
\[
X \sim \operatorname{Bin}(100,0.02)
\]

``Lieferung wird abgelehnt'' bei mehr als $4$ Ausschuss:
\[
P(X>4)=1-P(X\le 4)
\]

\textbf{3. Bitfehler}

Bei einem Byte:
\[
n=8
\]

Wenn ein Bit mit Wahrscheinlichkeit $p=0.2$ falsch ist:
\[
X \sim \operatorname{Bin}(8,0.2)
\]

\begin{align*}
P(\text{Byte korrekt}) &= P(X=0)=0.8^8 \\
P(X=2) &= \binom82 0.2^2 0.8^6 \\
P(X\le 2) &= \sum_{k=0}^2 \binom8k 0.2^k 0.8^{8-k}
\end{align*}

\subsection{Galton-Brett}
Beim Galton-Brett ist jede Ablenkung links/rechts ein Bernoulli-Experiment mit
\[
p=\frac12
\]

Nach $n$ Ebenen ist die Behaelternummer gleich der Anzahl der Rechtsablenkungen.

Also:
\[
X_n \sim \operatorname{Bin}\left(n,\frac12\right)
\]

\textbf{Merke:}
Galton-Brett = Binomialverteilung in Verkleidung.

\subsection{Geometrische Verteilung}
Hier wird nicht die Anzahl der Erfolge in fester Zeit gezaehlt, sondern:

\textbf{``Wie lange dauert es bis zum ersten Erfolg?''}

Wenn in jedem Versuch die Erfolgswahrscheinlichkeit $p$ ist und die Versuche unabhaengig sind, dann:
\[
X \sim \operatorname{Geom}(p)
\]

mit
\[
P(X=n)=p\,q^{n-1}
\qquad \text{fuer } n=1,2,3,\dots
\]

\textbf{Bedeutung:}
\begin{itemize}
  \item erst $n-1$ mal kein Erfolg
  \item dann einmal Erfolg
\end{itemize}

\subsection{Wann nehme ich die geometrische Verteilung?}
\begin{itemize}
  \item unbegrenzt viele Versuche moeglich
  \item jeder Versuch ist Bernoulli mit gleichem $p$
  \item gesucht ist die Wartezeit bis zum ersten Erfolg
\end{itemize}

\textbf{Typische Formulierungen}
\begin{itemize}
  \item ``beim wievielten Versuch erstmals Erfolg''
  \item ``wie lange dauert es bis zur ersten 6''
  \item ``mindestens so viele Versuche noetig''
\end{itemize}

\subsection{Wichtige geometrische Wahrscheinlichkeiten}
\[
P(X=n)=p\,q^{n-1}
\]

\[
P(X>n)=q^n
\]

\[
P(X\ge n)=q^{n-1}
\]

\[
P(X\le n)=1-q^n
\]

\textbf{Merke:}
``Noch kein Erfolg nach $n$ Versuchen'' bedeutet einfach:
\[
q^n
\]

\subsection{Typische geometrische Beispiele}
\textbf{1. Wuerfeln bis zur ersten 6}

\[
p=\frac16,\qquad q=\frac56
\]

\[
P(X=n)=\frac16\left(\frac56\right)^{n-1}
\]

\[
P(X>n)=\left(\frac56\right)^n
\]

\textbf{2. Mensch-aergere-dich-nicht}

Schon $20$ mal keine $6$ gewuerfelt?
Dann ist die Wahrscheinlichkeit fuer eine $6$ beim naechsten Wurf trotzdem wieder
\[
\frac16
\]

\textbf{Merke:}
Die geometrische Verteilung fuehrt genau zur Idee ``Vergangenheit egal'' bei unabhaengigen Versuchen. Vorheriges Pech macht den naechsten Erfolg nicht wahrscheinlicher.

\subsection{Binomial vs. geometrisch}
\begin{itemize}
  \item \textbf{Binomial:} Wie oft tritt Erfolg in $n$ Versuchen auf?
  \item \textbf{Geometrisch:} Wie viele Versuche braucht man bis zum ersten Erfolg?
\end{itemize}

\textbf{Kurzmerksatz}
\begin{itemize}
  \item \textbf{Anzahl von Erfolgen} $\Rightarrow$ binomial
  \item \textbf{Wartezeit bis erster Erfolg} $\Rightarrow$ geometrisch
\end{itemize}

\subsection{Hypergeometrische Verteilung}
Situation:
\begin{itemize}
  \item insgesamt $N$ Objekte
  \item davon $M$ ``Trefferobjekte'' / schwarze Kugeln
  \item es werden $n$ Objekte \textbf{ohne Zuruecklegen} gezogen
  \item $X$ = Anzahl Treffer unter den $n$ gezogenen
\end{itemize}

Dann gilt:
\[
X \sim \operatorname{Hyp}(N,M,n)
\]

mit
\[
P(X=k)=\frac{\binom{M}{k}\binom{N-M}{n-k}}{\binom{N}{n}}
\]

falls die Zahlen sinnvoll sind, sonst $0$.

\textbf{Bedeutung der Faktoren}
\begin{itemize}
  \item $\binom{M}{k}$ = waehle $k$ Trefferobjekte
  \item $\binom{N-M}{n-k}$ = waehle $n-k$ Nicht-Treffer
  \item $\binom{N}{n}$ = alle moeglichen Auswahlen
\end{itemize}

\subsection{Wann nehme ich die hypergeometrische Verteilung?}
\begin{itemize}
  \item feste Grundgesamtheit mit $N$ Objekten
  \item davon $M$ ``besondere''
  \item Ziehen \textbf{ohne Zuruecklegen}
  \item gesucht: Anzahl der besonderen Objekte in der Stichprobe
\end{itemize}

\textbf{Typische Formulierungen}
\begin{itemize}
  \item ``aus $100$ Dosen sind $10$ verunreinigt, $5$ werden entnommen''
  \item ``unter $20$ Losen sind $2$ Gewinne, $3$ Ziehungen ohne Zuruecklegen''
  \item ``Urnenziehen ohne Zuruecklegen''
\end{itemize}

\subsection{Typische hypergeometrische Beispiele}
\textbf{1. Erbsendosen / verunreinigte Dosen}

\[
N=100,\qquad M=10,\qquad n=5
\]

\[
P(X=k)=\frac{\binom{10}{k}\binom{90}{5-k}}{\binom{100}{5}}
\]

\[
P(X=0)=\frac{\binom{10}{0}\binom{90}{5}}{\binom{100}{5}}
\]

\[
P(X\le 1)=P(X=0)+P(X=1)
\]

\textbf{2. Lose ziehen ohne Zuruecklegen}

\[
N=20,\qquad M=2,\qquad n=3
\]

\[
P(\text{kein Gewinn})=P(X=0)=\frac{\binom20\binom{18}{3}}{\binom{20}{3}}
\]

\[
P(\text{genau ein Gewinn})=P(X=1)=\frac{\binom21\binom{18}{2}}{\binom{20}{3}}
\]

\subsection{Hypergeometrisch vs. binomial}
\textbf{Hypergeometrisch:}
\begin{itemize}
  \item ohne Zuruecklegen
  \item Trefferwahrscheinlichkeit aendert sich von Zug zu Zug
  \item Ziehungen sind nicht unabhaengig
\end{itemize}

\textbf{Binomial:}
\begin{itemize}
  \item mit Zuruecklegen oder modelliert als unabhaengig
  \item Trefferwahrscheinlichkeit bleibt konstant
  \item Ziehungen sind unabhaengig
\end{itemize}

\textbf{Wichtiger Grenzfall}
Wenn $n$ klein gegenueber $N$ ist, dann ist die hypergeometrische Verteilung oft sehr nah an der Binomialverteilung mit
\[
p \approx \frac{M}{N}
\]

\textbf{Merke:}
``Ohne Zuruecklegen'' schreit eigentlich nach hypergeometrisch. Binomial ist dann nur eine Annaeherung, wenn die Gesamtmenge sehr gross ist.

\subsection{Kurzvergleich der drei wichtigsten Verteilungen}
\[
\begin{array}{c|c|c}
\text{Verteilung} & \text{Frage} & \text{Formel} \\
\hline
\text{Binomial} & \text{Wie viele Erfolge in } n \text{ Versuchen?} &
\binom{n}{k}p^k q^{n-k} \\
\text{Geometrisch} & \text{Wann kommt der erste Erfolg?} &
p q^{n-1} \\
\text{Hypergeometrisch} & \text{Wie viele Treffer ohne Zuruecklegen?} &
\dfrac{\binom{M}{k}\binom{N-M}{n-k}}{\binom{N}{n}}
\end{array}
\]

\subsection{Typische Klausur-Fallen}
\begin{itemize}
  \item Binomial und hypergeometrisch verwechselt.
  \item Bei ``ohne Zuruecklegen'' trotzdem konstantes $p$ benutzt.
  \item Bei geometrischer Verteilung mit $n=0$ statt $n=1$ angefangen.
  \item ``mindestens einmal Erfolg'' nicht ueber das Gegenereignis gerechnet.
  \item Bei Binomial den Faktor $\binom{n}{k}$ vergessen.
  \item Bei geometrisch ``genau beim $n$-ten Versuch'' mit ``nach $n$ Versuchen'' verwechselt.
\end{itemize}

\subsection{Notation fuer saubere Loesungen}
Gute Standardsaetze:
\begin{itemize}
  \item ``Es handelt sich um $n$ unabhaengige Bernoulli-Versuche mit Erfolgswahrscheinlichkeit $p$. ''
  \item ``Damit ist $X$ binomialverteilt mit Parametern $n$ und $p$. ''
  \item ``Da ohne Zuruecklegen gezogen wird, ist $X$ hypergeometrisch verteilt. ''
  \item ``Gesucht ist die Wartezeit bis zum ersten Erfolg; daher verwenden wir die geometrische Verteilung. ''
  \item ``Wir berechnen die Wahrscheinlichkeit ueber das Gegenereignis. ''
\end{itemize}

\subsection{Mini-Merkliste}
\begin{itemize}
  \item feste Anzahl + unabhaengig + konstantes $p$ $\Rightarrow$ binomial
  \item erster Erfolg / Wartezeit $\Rightarrow$ geometrisch
  \item ohne Zuruecklegen $\Rightarrow$ hypergeometrisch
  \item Binomial: $\displaystyle P(X=k)=\binom{n}{k}p^k q^{n-k}$
  \item Geometrisch: $\displaystyle P(X=n)=p q^{n-1}$
  \item Hypergeometrisch: $\displaystyle P(X=k)=\frac{\binom{M}{k}\binom{N-M}{n-k}}{\binom{N}{n}}$
  \item ``mindestens einmal'' oft am schnellsten ueber $1-P(X=0)$ oder $1-q^n$
\end{itemize}

\section{Kapitel 1.7: Erwartungswert}

\subsection{Grundidee}
Der Erwartungswert ist der \textbf{gewichtete Mittelwert} einer Zufallsvariable.

Jeder moegliche Wert wird mit seiner Wahrscheinlichkeit gewichtet.

\textbf{Interpretation:}
\begin{itemize}
  \item langfristiger Durchschnitt bei sehr vielen Wiederholungen
  \item theoretischer Mittelwert
  \item nicht zwingend ein tatsaechlich moeglicher Einzelwert
\end{itemize}

\textbf{Wichtig:}
Ein Erwartungswert von $3.5$ beim Wuerfel heisst nicht, dass je eine $3.5$ fallen kann.

\subsection{Definition bei endlichem Wertebereich}
Hat $X$ die moeglichen Werte $x_1,\dots,x_n$, dann:
\[
E(X)=\sum_{i=1}^n x_i \cdot P(X=x_i)
\]

\subsection{Definition bei abz\"ahlbar unendlichem Wertebereich}
Hat $X$ die moeglichen Werte $x_1,x_2,\dots$, dann:
\[
E(X)=\sum_{i=1}^\infty x_i \cdot P(X=x_i)
\]

\textbf{Aber nur wenn die Reihe konvergiert.}

\textbf{Merke:}
Nicht jede Zufallsvariable hat automatisch einen Erwartungswert.

\subsection{Alternative Rechenform ueber $\Omega$}
Falls $p(\omega)=P(\{\omega\})$, dann gilt:
\[
E(X)=\sum_{\omega \in \Omega} X(\omega)\,p(\omega)
\]

\textbf{Wann nuetzlich?}
\begin{itemize}
  \item wenn $\Omega$ klein und konkret gegeben ist
  \item wenn $X(\omega)$ direkt ausgerechnet werden kann
  \item wenn mehrere $\omega$ denselben Funktionswert haben und man bequem gruppieren kann
\end{itemize}

\subsection{Standardrezept}
\begin{enumerate}
  \item Alle moeglichen Werte von $X$ bestimmen.
  \item Zu jedem Wert die Wahrscheinlichkeit bestimmen.
  \item Wert mal Wahrscheinlichkeit rechnen.
  \item Alles aufsummieren.
\end{enumerate}

\subsection{Wichtigste Rechenregeln}
\textbf{1. Linearitaet}
\begin{align*}
E(aX) &= aE(X) \\
E(X_1+X_2) &= E(X_1)+E(X_2)
\end{align*}

allgemeiner:
\[
E(a_1X_1+\dots+a_nX_n)=a_1E(X_1)+\dots+a_nE(X_n)
\]

\textbf{Extrem wichtig:}
Fuer die Summe brauchst du \textbf{keine} Unabhaengigkeit.

\subsection{2. Monotonie}
Wenn fuer alle $\omega$ gilt:
\[
X_1(\omega)\le X_2(\omega)
\]
dann folgt:
\[
E(X_1)\le E(X_2)
\]

\subsection{3. Erwartungswert von $g(X)$}
Wenn $Y=g(X)$, dann:
\[
E(Y)=E(g(X))=\sum_{x \in X(\Omega)} g(x)\,P(X=x)
\]

\textbf{Merke:}
Sehr praktisch bei $X^2$, $|X|$, Auszahlungen, Kostenfunktionen usw.

\subsection{4. Produktregel bei Unabhaengigkeit}
Wenn $X_1,X_2$ unabhaengig sind, dann:
\[
E(X_1X_2)=E(X_1)\cdot E(X_2)
\]

allgemeiner fuer unabhaengige $X_1,\dots,X_n$:
\[
E(X_1\cdots X_n)=E(X_1)\cdots E(X_n)
\]

\textbf{Wichtig:}
Ohne Unabhaengigkeit ist das im Allgemeinen \textbf{falsch}.

\subsection{Typische Klausurstrategie}
\textbf{Wenn eine Zufallsvariable eine Summe ist:}
\[
X=X_1+\dots+X_n
\]
dann fast immer sofort
\[
E(X)=E(X_1)+\dots+E(X_n)
\]
benutzen, statt alles brutal auszumultiplizieren.

\textbf{Wenn eine Zufallsvariable ein Produkt ist:}
\[
Y=X_1X_2
\]
dann nur bei Unabhaengigkeit sofort
\[
E(Y)=E(X_1)E(X_2)
\]

\subsection{Typische Beispiele}
\textbf{1. Fairer Wuerfel}
\[
E(X)=1\cdot \frac16+2\cdot \frac16+\dots+6\cdot \frac16=\frac{1+2+3+4+5+6}{6}=3.5
\]

\textbf{Merke:}
Bei Laplace-Verteilung ist der Erwartungswert einfach das arithmetische Mittel aller moeglichen Werte.

\subsection{2. Augensumme zweier fairer Wuerfel}
Seien
\[
X_1=\text{Augenzahl des ersten Wuerfels},\qquad X_2=\text{Augenzahl des zweiten Wuerfels}
\]

Dann:
\[
X=X_1+X_2
\]

und damit:
\[
E(X)=E(X_1)+E(X_2)=3.5+3.5=7
\]

\subsection{3. Augenprodukt zweier fairer Wuerfel}
\[
Y=X_1X_2
\]

Da $X_1$ und $X_2$ unabhaengig sind:
\[
E(Y)=E(X_1X_2)=E(X_1)E(X_2)=3.5^2=12.25
\]

\subsection{4. Bernoulli-Verteilung}
Wenn
\[
P(X=1)=p,\qquad P(X=0)=q=1-p
\]
dann:
\[
E(X)=0\cdot q+1\cdot p=p
\]

\textbf{Merke:}
Bernoulli-Mittelwert = Erfolgswahrscheinlichkeit.

\subsection{5. Binomialverteilung}
Wenn
\[
X \sim \operatorname{Bin}(n,p)
\]
dann:
\[
E(X)=np
\]

\textbf{Warum schnell?}
Man schreibt
\[
X=X_1+\dots+X_n
\]
mit Bernoulli-Variablen $X_i$, also:
\[
E(X)=E(X_1)+\dots+E(X_n)=p+\dots+p=np
\]

\subsection{6. Geometrische Verteilung}
Wenn
\[
X \sim \operatorname{Geom}(p)
\]
dann:
\[
E(X)=\frac{1}{p}
\]

\textbf{Beispiel}
Beim Warten auf die erste $6$:
\[
p=\frac16
\qquad \Rightarrow \qquad
E(X)=6
\]

\textbf{Merke:}
Geometrisch = mittlere Wartezeit bis zum ersten Erfolg.

\subsection{7. Hypergeometrische Verteilung}
Wenn
\[
X \sim \operatorname{Hyp}(N,M,n)
\]
dann:
\[
E(X)=n\cdot \frac{M}{N}
\]

\textbf{Interpretation:}
\[
\frac{M}{N}
\]
ist der Trefferanteil in der Grundgesamtheit, und davon zieht man im Mittel $n$-mal.

\subsection{8. Konstante Zufallsvariable}
Wenn
\[
X(\omega)=c \quad \text{fuer alle } \omega
\]
dann:
\[
E(X)=c
\]

\subsection{Gewinn-/Kostenaufgaben}
Wenn $X$ ein Gewinn in Cent oder Euro ist, dann ist
\[
E(X)
\]
der mittlere Gewinn pro Spiel.

\textbf{Interpretation}
\begin{itemize}
  \item $E(X)>0$: langfristig guenstig
  \item $E(X)<0$: langfristig schlecht
  \item $E(X)=0$: faires Spiel im Mittel
\end{itemize}

\textbf{Typische Klausuridee}
\begin{itemize}
  \item erst Auszahlungsvariable sauber definieren
  \item dann Erwartungswert ausrechnen
  \item bei Einsatz am Ende den Einsatz abziehen
\end{itemize}

\subsection{Wann nutze ich $g(X)$?}
Wenn nicht $X$ selbst, sondern z.B. das Quadrat, der Gewinn, die Kosten oder der Betrag gefragt ist:
\[
Y=g(X)
\]

dann:
\[
E(Y)=\sum_x g(x)\,P(X=x)
\]

\textbf{Beispiel}
Wenn $Y=X^2$, dann:
\[
E(Y)=E(X^2)=\sum_x x^2 P(X=x)
\]

\subsection{Typische Aufgabentypen}
\textbf{1. ``Wie viele im Mittel?''}

Das ist fast immer direkt eine Frage nach $E(X)$.

\textbf{2. Summe mehrerer Beitraege}

Dann fast immer Linearitaet benutzen.

\textbf{3. Produkt mehrerer Zufallsvariablen}

Nur bei Unabhaengigkeit auf Produkt der Erwartungswerte gehen.

\textbf{4. Warten bis erster Erfolg}

Dann geometrisch und
\[
E(X)=\frac1p
\]

\textbf{5. Anzahl Treffer in Stichprobe}

Je nach Situation:
\begin{itemize}
  \item mit Zuruecklegen / unabhaengig: binomial, $E(X)=np$
  \item ohne Zuruecklegen: hypergeometrisch, $E(X)=n\frac{M}{N}$
\end{itemize}

\subsection{Typische Fallen}
\begin{itemize}
  \item Erwartungswert mit wahrscheinlichstem Wert verwechselt.
  \item Erwartungswert als zwingend moeglichen Einzelwert interpretiert.
  \item Bei Summen unnoetig kompliziert gerechnet statt Linearitaet zu nutzen.
  \item Produktregel ohne Unabhaengigkeit benutzt.
  \item Bei $g(X)$ faelschlich $g(E(X))$ statt $E(g(X))$ gerechnet.
  \item Einsatz/Gewinn nicht sauber unterschieden.
  \item Bei unendlichem Wertebereich Konvergenz uebersehen.
\end{itemize}

\subsection{Notation fuer saubere Loesungen}
Gute Standardsaetze:
\begin{itemize}
  \item ``Der Erwartungswert berechnet sich als gewichtete Summe der moeglichen Werte.''
  \item ``Wir nutzen die Linearitaet des Erwartungswerts.''
  \item ``Da die Zufallsvariablen unabhaengig sind, gilt $E(X_1X_2)=E(X_1)E(X_2)$. ''
  \item ``Es handelt sich um eine geometrisch verteilte Zufallsvariable mit Parameter $p$, daher ist $E(X)=1/p$. ''
  \item ``Da $X$ binomialverteilt ist, gilt $E(X)=np$. ''
\end{itemize}

\subsection{Mini-Merkliste}
\begin{itemize}
  \item $\displaystyle E(X)=\sum_x x\,P(X=x)$
  \item alternativ: $\displaystyle E(X)=\sum_{\omega \in \Omega} X(\omega)\,p(\omega)$
  \item Linearitaet: $\displaystyle E(aX+bY)=aE(X)+bE(Y)$
  \item bei Unabhaengigkeit: $\displaystyle E(XY)=E(X)E(Y)$
  \item Bernoulli: $\displaystyle E(X)=p$
  \item Binomial: $\displaystyle E(X)=np$
  \item Geometrisch: $\displaystyle E(X)=1/p$
  \item Hypergeometrisch: $\displaystyle E(X)=n\frac{M}{N}$
\end{itemize}

\section{Kapitel 1.8: Varianz und Standardabweichung}

\subsection{Grundidee}
Der Erwartungswert sagt, \emph{wo} eine Zufallsvariable im Mittel liegt.
Die Varianz sagt, \emph{wie stark} sie um diesen Mittelwert streut.

\textbf{Merke:}
Gleicher Mittelwert heisst nicht gleiche Verteilung. Die Streuung kann komplett anders sein.

\subsection{Definition}
Sei
\[
\mu = E(X)
\]

Dann ist die Varianz:
\[
\operatorname{Var}(X)=D^2(X)=E\big((X-\mu)^2\big)
\]

Bei diskretem Wertebereich:
\[
\operatorname{Var}(X)=\sum_x (x-\mu)^2\,P(X=x)
\]

Die Standardabweichung ist:
\[
\sigma_X=\sqrt{\operatorname{Var}(X)}
\]

\textbf{Interpretation}
\begin{itemize}
  \item Varianz = mittlere quadratische Abweichung vom Mittelwert
  \item Standardabweichung = typische Streuung in derselben Einheit wie $X$
\end{itemize}

\subsection{Warum quadrieren?}
Weil sich sonst positive und negative Abweichungen wegheben wuerden.

\textbf{Merke:}
Varianz hat quadrierte Einheit, Standardabweichung wieder die normale Einheit.

\subsection{Alternative Rechenformel}
Die wichtigste Umformung:
\[
\operatorname{Var}(X)=E(X^2)-\big(E(X)\big)^2
\]

\textbf{Klausurregel:}
Fast immer zuerst mit dieser Formel rechnen, nicht direkt mit $(x-\mu)^2$ aufmachen, wenn $E(X)$ und $E(X^2)$ leicht zu bekommen sind.

\subsection{Rechenregeln}
\textbf{1. Skalierung und Verschiebung}
\begin{align*}
\operatorname{Var}(aX) &= a^2 \operatorname{Var}(X) \\
\operatorname{Var}(aX+b) &= a^2 \operatorname{Var}(X)
\end{align*}

\textbf{Merke:}
\begin{itemize}
  \item Multiplikation mit $a$ skaliert die Varianz mit $a^2$
  \item Addition einer Konstante $b$ aendert die Varianz gar nicht
\end{itemize}

\subsection{2. Summe zweier Zufallsvariablen}
Allgemein:
\[
\operatorname{Var}(X_1+X_2)
=\operatorname{Var}(X_1)+\operatorname{Var}(X_2)+2\big(E(X_1X_2)-E(X_1)E(X_2)\big)
\]

Das ist dasselbe wie:
\[
\operatorname{Var}(X_1+X_2)
=\operatorname{Var}(X_1)+\operatorname{Var}(X_2)+2\,\operatorname{cov}(X_1,X_2)
\]

Falls $X_1,X_2$ unabhaengig sind:
\[
\operatorname{Var}(X_1+X_2)=\operatorname{Var}(X_1)+\operatorname{Var}(X_2)
\]

\textbf{Extrem wichtig:}
Varianz ist \textbf{nicht} linear.

Also \textbf{falsch}:
\[
\operatorname{Var}(X_1+X_2)=\operatorname{Var}(X_1)+\operatorname{Var}(X_2)
\]
ohne Zusatzannahmen.

\subsection{3. Produkt bei Unabhaengigkeit}
Wenn $X_1,X_2$ unabhaengig sind:
\[
\operatorname{Var}(X_1X_2)=E(X_1^2)E(X_2^2)-\big(E(X_1)E(X_2)\big)^2
\]

\subsection{Spezialwerte wichtiger Verteilungen}
\textbf{Bernoulli-Verteilung}
\[
X \sim \operatorname{Bern}(p)
\]

\[
E(X)=p,\qquad \operatorname{Var}(X)=pq
\]
mit
\[
q=1-p
\]

\textbf{Binomialverteilung}
\[
X \sim \operatorname{Bin}(n,p)
\]

\[
E(X)=np,\qquad \operatorname{Var}(X)=npq
\]

\textbf{Geometrische Verteilung}
\[
X \sim \operatorname{Geom}(p)
\]

\[
E(X)=\frac1p,\qquad \operatorname{Var}(X)=\frac{q}{p^2}
\]

\textbf{Hypergeometrische Verteilung}
\[
X \sim \operatorname{Hyp}(N,M,n)
\]

\[
E(X)=n\frac{M}{N}
\]

\[
\operatorname{Var}(X)=
n\cdot \frac{M}{N}\left(1-\frac{M}{N}\right)\cdot \frac{N-n}{N-1}
\]

\textbf{Merke:}
Die hypergeometrische Varianz sieht fast aus wie bei der Binomialverteilung, aber mit dem Zusatzfaktor
\[
\frac{N-n}{N-1}
\]
wegen Ziehen ohne Zuruecklegen.

\subsection{Beispiele}
\textbf{1. Fairer Wuerfel}
\[
E(X)=3.5,\qquad E(X^2)=\frac{1^2+2^2+\dots+6^2}{6}=\frac{91}{6}
\]

also
\[
\operatorname{Var}(X)=\frac{91}{6}-3.5^2=\frac{35}{12}
\]

\textbf{2. Summe zweier fairer Wuerfel}

Mit Unabhaengigkeit:
\[
\operatorname{Var}(X_1+X_2)=\operatorname{Var}(X_1)+\operatorname{Var}(X_2)=2\cdot \frac{35}{12}=\frac{35}{6}
\]

\subsection{Varianz Null}
\[
\operatorname{Var}(X)=0
\]
genau dann, wenn $X$ fast sicher konstant ist.

\textbf{Merke:}
Keine Streuung $\Leftrightarrow$ immer derselbe Wert mit Wahrscheinlichkeit $1$.

\subsection{Tschebyscheff-Ungleichung}
Fuer jede Zufallsvariable mit Mittelwert $\mu$ und Varianz $\sigma^2$ gilt fuer jedes $c>0$:
\[
P(|X-\mu|\ge c)\le \frac{\sigma^2}{c^2}
\]

und damit auch:
\[
P(|X-\mu|<c)\ge 1-\frac{\sigma^2}{c^2}
\]

\textbf{Bedeutung:}
Große Abweichungen vom Mittelwert koennen durch die Varianz nach oben abgeschaetzt werden.

\textbf{Wann benutzen?}
\begin{itemize}
  \item wenn nur Mittelwert und Varianz bekannt sind
  \item wenn eine grobe obere Schranke fuer Abweichungswahrscheinlichkeiten gesucht ist
\end{itemize}

\subsection{Typische Klausuraufgaben}
\textbf{1. Varianz direkt berechnen}

Rezept:
\begin{enumerate}
  \item $E(X)$ berechnen
  \item $E(X^2)$ berechnen
  \item $\operatorname{Var}(X)=E(X^2)-E(X)^2$
\end{enumerate}

\textbf{2. Varianz von Summen}

Wenn unabhaengig:
\[
\operatorname{Var}(X_1+\dots+X_n)=\operatorname{Var}(X_1)+\dots+\operatorname{Var}(X_n)
\]

\textbf{3. Standardabweichung}

Einfach:
\[
\sigma=\sqrt{\operatorname{Var}(X)}
\]

\textbf{4. Abschaetzungen mit Tschebyscheff}

Wenn gefragt wird ``mit welcher Wahrscheinlichkeit liegt $X$ hoechstens $c$ vom Mittelwert entfernt?'', dann:
\[
P(|X-\mu|<c)\ge 1-\frac{\sigma^2}{c^2}
\]

\subsection{Typische Fallen}
\begin{itemize}
  \item Varianz mit Standardabweichung verwechselt.
  \item Bei $\operatorname{Var}(aX+b)$ das $a^2$ vergessen.
  \item Bei Summen die Unabhaengigkeit stillschweigend angenommen.
  \item $E(X^2)$ mit $(E(X))^2$ verwechselt.
  \item Tschebyscheff als exakte Wahrscheinlichkeit statt als Schranke gelesen.
\end{itemize}

\subsection{Mini-Merkliste}
\begin{itemize}
  \item $\displaystyle \operatorname{Var}(X)=E((X-\mu)^2)$
  \item $\displaystyle \operatorname{Var}(X)=E(X^2)-E(X)^2$
  \item $\displaystyle \sigma_X=\sqrt{\operatorname{Var}(X)}$
  \item $\displaystyle \operatorname{Var}(aX+b)=a^2\operatorname{Var}(X)$
  \item unabhaengig: $\displaystyle \operatorname{Var}(X+Y)=\operatorname{Var}(X)+\operatorname{Var}(Y)$
  \item Bernoulli: $\displaystyle pq$
  \item Binomial: $\displaystyle npq$
  \item Geometrisch: $\displaystyle q/p^2$
  \item Hypergeometrisch: $\displaystyle n\frac{M}{N}(1-\frac{M}{N})\frac{N-n}{N-1}$
\end{itemize}

\section{Kapitel 1.9: Kovarianz und Korrelation}

\subsection{Grundidee}
Jetzt geht es nicht mehr um die Streuung \emph{einer} Variablen, sondern um den Zusammenhang \emph{zweier} Variablen.

\textbf{Fragen}
\begin{itemize}
  \item Gehen große Werte von $X$ oft mit großen Werten von $Y$ einher?
  \item Oder eher mit kleinen Werten?
  \item Wie stark ist dieser Zusammenhang?
\end{itemize}

\subsection{Kovarianz}
Die Kovarianz ist definiert als
\[
\operatorname{cov}(X,Y)=E\big((X-\mu_X)(Y-\mu_Y)\big)
\]

\textbf{Interpretation des Vorzeichens}
\begin{itemize}
  \item $\operatorname{cov}(X,Y)>0$: eher gleichlaufend
  \item $\operatorname{cov}(X,Y)<0$: eher gegenlaeufig
  \item $\operatorname{cov}(X,Y)=0$: kein linearer Zusammenhang messbar
\end{itemize}

\textbf{Wichtig:}
Korrelation bedeutet \textbf{nicht} Kausalitaet.

\subsection{Wichtigste Umformung}
\[
\operatorname{cov}(X,Y)=E(XY)-E(X)E(Y)
\]

\textbf{Klausurregel:}
Fast immer mit dieser Formel rechnen, nicht direkt mit $(X-\mu_X)(Y-\mu_Y)$ aufmachen.

\subsection{Wichtige Regeln}
\textbf{1. Bilinearitaet}
\begin{align*}
\operatorname{cov}(X_1+X_2,Y) &= \operatorname{cov}(X_1,Y)+\operatorname{cov}(X_2,Y) \\
\operatorname{cov}(aX,Y) &= a\,\operatorname{cov}(X,Y)
\end{align*}

analog auch im zweiten Argument.

\textbf{2. Symmetrie}
\[
\operatorname{cov}(X,Y)=\operatorname{cov}(Y,X)
\]

\textbf{3. Spezialfall Varianz}
\[
\operatorname{cov}(X,X)=\operatorname{Var}(X)
\]

\subsection{Unabhaengigkeit und Kovarianz}
Wenn $X,Y$ unabhaengig sind, dann:
\[
\operatorname{cov}(X,Y)=0
\]

\textbf{Aber Achtung:}
Die Umkehrung gilt im Allgemeinen \textbf{nicht}.

\textbf{Merke:}
Unabhaengig $\Rightarrow$ unkorreliert, aber unkorreliert $\nRightarrow$ unabhaengig.

\subsection{Varianz von Summen mit Kovarianz}
\[
\operatorname{Var}(X+Y)=\operatorname{Var}(X)+\operatorname{Var}(Y)+2\operatorname{cov}(X,Y)
\]

Allgemeiner:
\[
\operatorname{Var}(X_1+\dots+X_n)
=\sum_{i=1}^n \operatorname{Var}(X_i)
+2\sum_{1\le i<j\le n}\operatorname{cov}(X_i,X_j)
\]

Falls die $X_i$ paarweise unabhaengig sind:
\[
\operatorname{Var}(X_1+\dots+X_n)=\sum_{i=1}^n \operatorname{Var}(X_i)
\]

\subsection{Beispiel: Wuerfelsumme}
Bei $n$ unabhaengigen fairen Wuerfeln:
\[
X=X_1+\dots+X_n
\]

\[
E(X)=3.5n
\]

und wegen Unabhaengigkeit:
\[
\operatorname{Var}(X)=n\cdot \frac{35}{12}
\]

\subsection{Korrelation / Korrelationskoeffizient}
Die Kovarianz haengt von der Skala ab. Deshalb normiert man:
\[
\rho(X,Y)=\frac{\operatorname{cov}(X,Y)}{\sigma_X\sigma_Y}
\]

vorausgesetzt $\sigma_X>0$ und $\sigma_Y>0$.

\textbf{Eigenschaften}
\[
-1 \le \rho(X,Y)\le 1
\]

\textbf{Interpretation}
\begin{itemize}
  \item $\rho>0$: gleichlaufender linearer Zusammenhang
  \item $\rho<0$: gegenlaeufiger linearer Zusammenhang
  \item $|\rho|$ groß: starker linearer Zusammenhang
  \item $\rho=0$: kein linearer Zusammenhang messbar
\end{itemize}

\subsection{Extremfaelle}
\[
\rho(X,Y)=1
\]
bedeutet: $Y$ ist fast sicher eine positiv lineare Funktion von $X$.

\[
\rho(X,Y)=-1
\]
bedeutet: $Y$ ist fast sicher eine negativ lineare Funktion von $X$.

\textbf{Merke:}
$\rho=0$ heisst nicht automatisch unabhaengig, sondern nur: kein \emph{linearer} Zusammenhang.

\subsection{Standardisierung}
Die standardisierte Variable ist:
\[
X^*=\frac{X-\mu_X}{\sigma_X}
\]

Dann gilt immer:
\[
E(X^*)=0,\qquad \operatorname{Var}(X^*)=1
\]

und
\[
\rho(X,Y)=\operatorname{cov}(X^*,Y^*)
\]

\subsection{Typische Klausuraufgaben}
\textbf{1. Kovarianz berechnen}

Rezept:
\begin{enumerate}
  \item $E(X)$ berechnen
  \item $E(Y)$ berechnen
  \item $E(XY)$ berechnen
  \item $\operatorname{cov}(X,Y)=E(XY)-E(X)E(Y)$
\end{enumerate}

\textbf{2. Korrelation berechnen}

Danach:
\[
\rho(X,Y)=\frac{\operatorname{cov}(X,Y)}{\sigma_X\sigma_Y}
\]

\textbf{3. Unabhaengigkeit pruefen}

Nicht aus $\rho=0$ schliessen.
Lieber direkt ueber gemeinsame Verteilung oder Definition pruefen.

\subsection{Typische Fallen}
\begin{itemize}
  \item Korrelation mit Kausalitaet verwechselt.
  \item Aus Kovarianz $0$ auf Unabhaengigkeit geschlossen.
  \item Normierung mit Standardabweichung vergessen.
  \item Bei Summen die Kovarianzterme vergessen.
  \item $\rho=\pm 1$ falsch interpretiert.
\end{itemize}

\subsection{Mini-Merkliste}
\begin{itemize}
  \item $\displaystyle \operatorname{cov}(X,Y)=E((X-\mu_X)(Y-\mu_Y))$
  \item $\displaystyle \operatorname{cov}(X,Y)=E(XY)-E(X)E(Y)$
  \item unabhaengig $\Rightarrow \operatorname{cov}(X,Y)=0$
  \item $\displaystyle \operatorname{Var}(X+Y)=\operatorname{Var}(X)+\operatorname{Var}(Y)+2\operatorname{cov}(X,Y)$
  \item $\displaystyle \rho(X,Y)=\frac{\operatorname{cov}(X,Y)}{\sigma_X\sigma_Y}$
  \item $\displaystyle -1\le \rho(X,Y)\le 1$
  \item $\rho=0$ bedeutet nur: unkorreliert, nicht automatisch unabhaengig
\end{itemize}

\section{Kapitel 1.10: Schwaches Gesetz der grossen Zahlen}

\subsection{Grundidee}
Wenn man ein Zufallsexperiment oft genug unabhaengig wiederholt, dann liegt das arithmetische Mittel der beobachteten Werte mit hoher Wahrscheinlichkeit nahe am Erwartungswert.

\textbf{Anschaulich:}
Viele Wiederholungen stabilisieren den Durchschnitt.

\subsection{Situation}
Sei $X$ eine Zufallsvariable mit
\[
E(X)=\mu,\qquad \operatorname{Var}(X)=\sigma^2
\]

Seien $X_1,\dots,X_n$ unabhaengige Realisierungen derselben Verteilung.

Dann definiert man das Stichprobenmittel:
\[
Z_n=\frac1n\sum_{i=1}^n X_i
\]

\subsection{Erwartungswert und Varianz des Mittels}
\[
E(Z_n)=\mu
\]

\[
\operatorname{Var}(Z_n)=\frac{\sigma^2}{n}
\]

\textbf{Merke:}
Der Mittelwert bleibt gleich, aber die Streuung wird mit wachsendem $n$ kleiner.

\subsection{Schwaches Gesetz der grossen Zahlen}
Fuer jedes $c>0$ gilt:
\[
P(|Z_n-\mu|\ge c)\le \frac{\sigma^2}{n c^2}
\]

und daher
\[
\lim_{n\to\infty} P(|Z_n-\mu|\ge c)=0
\]

\textbf{Bedeutung:}
Fuer jede feste Toleranz $c$ wird die Wahrscheinlichkeit einer groesseren Abweichung beliebig klein, wenn $n$ gross genug ist.

\textbf{Merke:}
Das ist die praezise mathematische Form von
\[
\text{``Durchschnitt } \approx \text{ Erwartungswert bei vielen Wiederholungen.''}
\]

\subsection{Woher kommt das?}
Im Wesentlichen aus:
\begin{itemize}
  \item Linearitaet des Erwartungswerts
  \item Varianz von Summen unabhaengiger Variablen
  \item Tschebyscheff-Ungleichung
\end{itemize}

\subsection{Relative Haeufigkeit als Spezialfall}
Sei $A$ ein Ereignis und $X_A$ seine Indikatorvariable:
\[
X_A(\omega)=
\begin{cases}
1,& \omega \in A \\
0,& \omega \notin A
\end{cases}
\]

Dann gilt:
\[
E(X_A)=P(A)=p
\]

und
\[
\operatorname{Var}(X_A)=p(1-p)=pq
\]

Bei $n$ Wiederholungen ist die relative Haeufigkeit
\[
r_n(A)=\frac{n_A}{n}
\]
gerade das arithmetische Mittel der Indikatorvariablen.

\textbf{Also:}
\[
r_n(A)\approx p
\]
fuer grosses $n$.

\subsection{Praezise Abschaetzung fuer relative Haeufigkeiten}
Fuer jedes $c>0$ gilt:
\[
P(|r_n(A)-p|\ge c)\le \frac{pq}{n c^2}
\]

Da immer
\[
pq\le \frac14
\]
gilt, folgt die grobere, aber praktische Schranke:
\[
P(|r_n(A)-p|\ge c)\le \frac{1}{4nc^2}
\]

\textbf{Umgestellt:}
\[
P(|r_n(A)-p|<c)\ge 1-\frac{1}{4nc^2}
\]

\subsection{Wofuer ist das gut?}
Damit kann man abschaetzen:
\begin{itemize}
  \item wie viele Wiederholungen noetig sind
  \item um eine Wahrscheinlichkeit $p$ auf Genauigkeit $c$ zu schaetzen
  \item mit vorgegebener Sicherheit
\end{itemize}

\subsection{Typische Klausurform}
Gesucht ist oft ein moeglichst kleines $n$, so dass
\[
P(|r_n(A)-p|<c)\ge \alpha
\]

Mit der groben Schranke reicht:
\[
1-\frac{1}{4nc^2}\ge \alpha
\]

also
\[
\frac{1}{4nc^2}\le 1-\alpha
\]

und daraus $n$ ausrechnen.

\subsection{Beispieltyp}
``Mit mindestens $90\%$ Sicherheit soll die Abweichung hoechstens $0.05$ sein.''

Dann:
\[
1-\frac{1}{4n\cdot 0.05^2}\ge 0.9
\]

und daraus folgt das benoetigte $n$.

\subsection{Interpretation fuer Messungen}
Wenn man eine physikalische Groesse mehrfach unabhaengig misst:
\begin{itemize}
  \item Einzelmessungen streuen
  \item der Mittelwert vieler Messungen ist stabiler
  \item die Unsicherheit des Mittels sinkt mit wachsendem $n$
\end{itemize}

\textbf{Merke:}
Mehrfach messen bringt nicht magisch den exakten Wert, aber das arithmetische Mittel wird verlaesslicher.

\subsection{Typische Fallen}
\begin{itemize}
  \item Schwaches Gesetz mit ``sicherer Konvergenz'' verwechseln.
  \item Glauben, dass schon kleines $n$ automatisch gute Genauigkeit garantiert.
  \item Relative Haeufigkeit mit Wahrscheinlichkeit exakt gleichsetzen.
  \item $pq\le 1/4$ vergessen, wenn $p$ unbekannt ist.
\end{itemize}

\subsection{Mini-Merkliste}
\begin{itemize}
  \item $\displaystyle Z_n=\frac1n\sum_{i=1}^n X_i$
  \item $\displaystyle E(Z_n)=\mu$
  \item $\displaystyle \operatorname{Var}(Z_n)=\frac{\sigma^2}{n}$
  \item $\displaystyle P(|Z_n-\mu|\ge c)\le \frac{\sigma^2}{nc^2}$
  \item relative Haeufigkeit ist Mittelwert von Indikatorvariablen
  \item $\displaystyle P(|r_n(A)-p|\ge c)\le \frac{pq}{nc^2}\le \frac{1}{4nc^2}$
\end{itemize}

\section{Kapitel 1.11: Anwendungsbeispiele}

\subsection{Worum es hier eigentlich geht}
Das Kapitel bringt kaum neue Theorie, sondern zeigt, wie man das bisherige Zeug auf reale Modellierungen wirft.

\textbf{Wichtig fuer die Klausur:}
Nicht neue Formeln lernen, sondern erkennen:
\begin{itemize}
  \item Was ist die Ergebnismenge?
  \item Was sind sinnvolle Ereignisse?
  \item Wo steckt bedingte Wahrscheinlichkeit?
  \item Wo braucht man totale Wahrscheinlichkeit / Bayes?
  \item Wo steckt Binomialverteilung / Erwartungswert / Tschebyscheff?
\end{itemize}

\subsection{Beispiel Uebertragungskanal}
Typische Modellierung:
\begin{itemize}
  \item Senderalphabet:
  \[
  S=\{s_1,\dots,s_m\}
  \]
  \item Empfaengeralphabet:
  \[
  E=\{e_1,\dots,e_n\}
  \]
  \item Ergebnismenge:
  \[
  \Omega=S\times E
  \]
\end{itemize}

Ereignisse:
\begin{align*}
S_i &= \{\text{$s_i$ wurde gesendet}\} \\
E_j &= \{\text{$e_j$ wurde empfangen}\}
\end{align*}

\subsection{Wichtige Groessen beim Kanal}
\textbf{1. A-priori-Wahrscheinlichkeiten}
\[
p_i=P(S_i)
\]

\textbf{2. Kanalwahrscheinlichkeiten}
\[
p_{ij}=P(E_j\mid S_i)
\]

Das ist eine Matrix bedingter Wahrscheinlichkeiten.

\textbf{Wichtige Bedingung}
Fuer jedes feste $i$:
\[
\sum_{j=1}^n p_{ij}=1
\]

\textbf{Merke:}
Eine Zeile = gesendetes Zeichen fest, dann muessen sich die Empfangswahrscheinlichkeiten zu $1$ summieren.

\subsection{Gemeinsame Wahrscheinlichkeiten}
Mit Multiplikationssatz:
\[
P(S_i \cap E_j)=P(E_j\mid S_i)\,P(S_i)=p_{ij}p_i
\]

\subsection{Empfangswahrscheinlichkeiten}
Mit totaler Wahrscheinlichkeit:
\[
P(E_j)=\sum_{i=1}^m P(E_j\mid S_i)P(S_i)=\sum_{i=1}^m p_{ij}p_i
\]

\subsection{Rueckschluss auf gesendetes Zeichen}
Mit Bayes:
\[
P(S_i\mid E_j)=\frac{P(E_j\mid S_i)P(S_i)}{P(E_j)}
=\frac{p_{ij}p_i}{\sum_{k=1}^m p_{kj}p_k}
\]

\textbf{Klausurregel:}
Wenn gefragt wird ``welches Zeichen wurde wahrscheinlich gesendet, wenn $e_j$ empfangen wurde?'', dann Bayes.

\subsection{Entscheidungsfunktion}
Eine Entscheidungsfunktion
\[
\varphi(j)
\]
sagt:
``Wenn $e_j$ empfangen wird, vermuten wir $s_{\varphi(j)}$ als gesendet.''

\textbf{Bestes Prinzip}
Waehle fuer jedes $j$ das $i$, fuer das
\[
P(S_i\mid E_j)
\]
maximal ist.

\textbf{Merke:}
Bei gegebener Empfangsbeobachtung nimmt man den wahrscheinlichsten Senderzustand.

\subsection{Fehlerwahrscheinlichkeit}
Wenn $K$ = ``korrekt entschieden'', dann:
\[
P(K)=\sum_j P(S_{\varphi(j)}\cap E_j)
\]

und
\[
P(\text{Fehler})=1-P(K)
\]

\subsection{Viele Uebertragungen}
Bei $n$ unabhaengigen Uebertragungen und Fehlerwahrscheinlichkeit $p_f$ pro Zeichen:
\[
X=\text{Anzahl Fehler}
\sim \operatorname{Bin}(n,p_f)
\]

Dann:
\[
E(X)=np_f,\qquad \operatorname{Var}(X)=np_f(1-p_f)
\]

und fuer Abschaetzungen grosser Abweichungen kann Tschebyscheff benutzt werden.

\subsection{Mini-Merkliste}
\begin{itemize}
  \item Kanalmatrix = bedingte Wahrscheinlichkeiten $P(E_j\mid S_i)$
  \item gemeinsame Wkt.: Multiplikationssatz
  \item Empfangswkt.: totale Wahrscheinlichkeit
  \item Rueckschluss Sender aus Empfang: Bayes
  \item viele Fehler in Folge: Binomialverteilung
\end{itemize}

\section{Kapitel 2.1: Allgemeine Wahrscheinlichkeitsraeume}

\subsection{Warum reicht der diskrete Fall nicht mehr?}
In Kapitel 1 war $\Omega$ endlich oder abz\"ahlbar unendlich.

Das reicht nicht fuer Groessen wie:
\begin{itemize}
  \item Lebensdauer eines Atomkerns
  \item Abfuellmenge einer Flasche
  \item Messfehler
  \item beliebige reelle Messwerte
\end{itemize}

Da liegen die moeglichen Werte oft in Intervallen wie
\[
[0,\infty),\qquad \mathbb{R},\qquad [a,b]
\]
und das ist ueberabzaehlbar.

\subsection{Was aendert sich?}
Im allgemeinen Fall nimmt man nicht mehr einfach alle Teilmengen von $\Omega$ als Ereignisse.

Stattdessen:
\begin{itemize}
  \item $\Omega$ = Ergebnismenge
  \item $\mathcal A$ = Menge der zulaessigen Ereignisse
  \item $P:\mathcal A \to [0,1]$ = Wahrscheinlichkeitsmass
\end{itemize}

Der allgemeine Wahrscheinlichkeitsraum ist also:
\[
(\Omega,\mathcal A,P)
\]

\subsection{Sigma-Algebra}
$\mathcal A$ heisst Ereignis-$\sigma$-Algebra.

Fuer die Klausur reicht im Grunde:
\begin{itemize}
  \item $\Omega \in \mathcal A$
  \item $\emptyset \in \mathcal A$
  \item mit einem Ereignis auch sein Komplement
  \item mit abz\"ahlbar vielen Ereignissen auch deren Vereinigung und Durchschnitt
\end{itemize}

\textbf{Merke:}
Eine $\sigma$-Algebra ist einfach ein Ereignissystem, mit dem man sauber weiterrechnen kann.

\subsection{Borel-Sigma-Algebra}
Fuer Zufallsvariablen mit Werten in $\mathbb R$ nimmt man typischerweise die Borel-$\sigma$-Algebra
\[
\mathcal B(\mathbb R)
\]

Die enthaelt insbesondere:
\begin{itemize}
  \item alle Intervalle
  \item alle einpunktigen Mengen $\{x\}$
  \item alle endlichen und abz\"ahlbaren Teilmengen von $\mathbb R$
\end{itemize}

\textbf{Merke:}
Praktisch relevante Ereignisse wie
\[
X\le b,\quad a<X\le b,\quad X\in [a,b]
\]
sind damit sauber zulaessig.

\subsection{Allgemeine Zufallsvariable}
Im allgemeinen Fall ist eine Zufallsvariable eine Abbildung
\[
X:\Omega\to \mathbb R
\]
mit der Eigenschaft, dass fuer jedes Borel-Ereignis $A \in \mathcal B(\mathbb R)$ gilt:
\[
X^{-1}(A)\in \mathcal A
\]

\textbf{Bedeutung:}
Man darf alle Ereignisse der Form
\[
\{X\in A\}
\]
wahrscheinlichkeitstheoretisch auswerten.

\subsection{Verteilung von $X$}
Wie vorher:
\[
P_X(A)=P(X\in A)
\qquad \text{fuer } A\in \mathcal B(\mathbb R)
\]

\textbf{Wichtig:}
Die Idee aus Kapitel 1 bleibt gleich. Nur die erlaubten Mengen werden sauberer formuliert.

\subsection{Diskret vs. allgemein}
\textbf{Diskrete Zufallsvariable}
\begin{itemize}
  \item Wertemenge endlich oder abz\"ahlbar
  \item Beschreibung oft ueber Wahrscheinlichkeitsfunktion $p(x)$
\end{itemize}

\textbf{Allgemeiner / stetiger Fall}
\begin{itemize}
  \item Wertemenge oft Intervall oder ganz $\mathbb R$
  \item Beschreibung spaeter ueber Verteilungsfunktion bzw. Dichte
\end{itemize}

\subsection{Unabhaengigkeit im allgemeinen Fall}
Zwei Zufallsvariablen $X,Y$ sind unabhaengig, wenn fuer alle Borel-Mengen $A,B$ gilt:
\[
P(X\in A,\;Y\in B)=P(X\in A)\cdot P(Y\in B)
\]

Oft reicht es im stetigen Kontext zu denken an:
\[
P(X\le a,\;Y\le b)=P(X\le a)\cdot P(Y\le b)
\]

\subsection{Was bleibt gleich?}
Fast alle Konzepte aus Kapitel 1 bleiben inhaltlich gleich:
\begin{itemize}
  \item Ereignisse
  \item bedingte Wahrscheinlichkeit
  \item Unabhaengigkeit
  \item Erwartungswert
  \item Varianz
\end{itemize}

Nur die technische Sprache wird etwas praeziser.

\subsection{Mini-Merkliste}
\begin{itemize}
  \item allgemein: $\displaystyle (\Omega,\mathcal A,P)$
  \item nicht mehr jede Teilmenge muss Ereignis sein
  \item fuer reelle Werte wichtig: Borel-$\sigma$-Algebra $\mathcal B(\mathbb R)$
  \item Zufallsvariable muss Borel-Ereignisse auf Ereignisse in $\Omega$ zurueckfuehren
  \item Inhaltlich bleibt fast alles aus Kapitel 1 erhalten
\end{itemize}

\section{Kapitel 2.2: Dichte und Verteilungsfunktion}

\subsection{Verteilungsfunktion}
Zu einer Zufallsvariablen $X$ gehoert die Verteilungsfunktion
\[
F_X(x)=P(X\le x)
\]

Oft kurz:
\[
F(x)=P(X\le x)
\]

\textbf{Bedeutung:}
Sie sagt fuer jedes $x$, mit welcher Wahrscheinlichkeit $X$ hoechstens $x$ ist.

\subsection{Wahrscheinlichkeiten aus $F$ ablesen}
\[
P(a<X\le b)=F(b)-F(a)
\]

allgemeiner:
\begin{align*}
P(a<X<b) &= F(b^-)-F(a) \\
P(a\le X<b) &= F(b^-)-F(a^-) \\
P(a\le X\le b) &= F(b)-F(a^-)
\end{align*}

\textbf{Im stetigen Fall} vereinfacht sich das spaeter stark, weil dann Punktwahrscheinlichkeiten $0$ sind.

\subsection{Eigenschaften jeder Verteilungsfunktion}
Fuer jede Verteilungsfunktion gilt:
\begin{itemize}
  \item $0\le F(x)\le 1$
  \item $F$ ist monoton wachsend
  \item $\lim_{x\to -\infty}F(x)=0$
  \item $\lim_{x\to \infty}F(x)=1$
  \item $F$ ist rechtsstetig
\end{itemize}

\subsection{Spruenge und Punktwahrscheinlichkeiten}
\[
P(X=x)=F(x)-F(x^-)
\]

\textbf{Merke:}
Die Sprunghoehe bei $x$ ist genau die Wahrscheinlichkeit der Stelle $x$.

Also:
\begin{itemize}
  \item Sprung vorhanden $\Rightarrow P(X=x)>0$
  \item keine Spruenge / stetig an $x$ $\Rightarrow P(X=x)=0$
\end{itemize}

\subsection{Stetige Zufallsvariable}
$X$ heisst stetig, wenn es eine Dichte $f$ gibt mit
\[
F(x)=\int_{-\infty}^x f(u)\,du
\]

Dann ist $f$ die Wahrscheinlichkeitsdichte von $X$.

\subsection{Wahrscheinlichkeiten mit Dichte}
Wenn $X$ Dichte $f$ hat, dann:
\[
P(X\le x)=F(x)=\int_{-\infty}^x f(u)\,du
\]

und
\[
P(a<X\le b)=\int_a^b f(u)\,du
\]

\textbf{Grafische Bedeutung:}
Wahrscheinlichkeiten sind Flaechen unter der Dichte.

\subsection{Wichtige Eigenschaften einer Dichte}
Eine Dichte $f$ muss erfuellen:
\[
f(x)\ge 0
\quad \text{(fast ueberall)}
\]

und
\[
\int_{-\infty}^{\infty} f(x)\,dx=1
\]

\textbf{Merke:}
Eine Dichte ist keine Wahrscheinlichkeit an einer Stelle, sondern eine Flaechendichte.

\subsection{Extrem wichtig}
Bei stetigen Zufallsvariablen gilt:
\[
P(X=x)=0
\quad \text{fuer alle } x
\]

Deshalb:
\[
P(a\le X\le b)=P(a<X\le b)=P(a<X<b)=P(a\le X<b)
\]

\textbf{Merke:}
Bei stetigen Variablen sind die Randklammern egal.

\subsection{Zusammenhang zwischen Dichte und Verteilungsfunktion}
Wenn $F$ differenzierbar ist, dann:
\[
f(x)=F'(x)
\]

und umgekehrt:
\[
F(x)=\int_{-\infty}^x f(u)\,du
\]

\textbf{Kurzmerksatz}
\begin{itemize}
  \item Dichte integrieren $\Rightarrow$ Verteilungsfunktion
  \item Verteilungsfunktion ableiten $\Rightarrow$ Dichte
\end{itemize}

\subsection{Stetige Gleichverteilung auf $[a,b]$}
Dichte:
\[
f(x)=
\begin{cases}
\frac{1}{b-a},& x\in [a,b] \\
0,& \text{sonst}
\end{cases}
\]

Verteilungsfunktion:
\[
F(x)=
\begin{cases}
0,& x<a \\
\frac{x-a}{b-a},& a\le x\le b \\
1,& x>b
\end{cases}
\]

\textbf{Merke:}
Bei Gleichverteilung auf $[a,b]$ ist Wahrscheinlichkeit proportional zur Intervalllaenge.

Also:
\[
P(c\le X\le d)=\frac{d-c}{b-a}
\qquad \text{fuer } a\le c\le d\le b
\]

\subsection{Exponentialverteilung}
Mit Parameter $\lambda>0$:

\[
f(x)=
\begin{cases}
0,& x<0 \\
\lambda e^{-\lambda x},& x\ge 0
\end{cases}
\]

\[
F(x)=
\begin{cases}
0,& x<0 \\
1-e^{-\lambda x},& x\ge 0
\end{cases}
\]

\textbf{Wichtig:}
\[
P(X\ge a)=1-F(a)=e^{-\lambda a}
\qquad (a\ge 0)
\]

\textbf{Merke:}
Exponentialverteilung modelliert Wartezeiten und ist das stetige Pendant zur geometrischen Verteilung.

\subsection{Gedaechtnislosigkeit}
Die Exponentialverteilung ist gedaechtnislos:
\[
P(X>n+m\mid X>n)=P(X>m)
\]

\textbf{Merke:}
Schon lange gewartet zu haben macht das baldige Eintreten nicht wahrscheinlicher.

\subsection{Diskret vs. stetig}
\textbf{Diskret}
\begin{itemize}
  \item Wahrscheinlichkeitsfunktion $p(x)$
  \item Wahrscheinlichkeiten durch Summen
\end{itemize}

\textbf{Stetig}
\begin{itemize}
  \item Dichte $f(x)$
  \item Wahrscheinlichkeiten durch Integrale
\end{itemize}

\textbf{Analogie}
\[
P(a\le X\le b)=\sum_{x\in[a,b]} p(x)
\qquad \text{vs.} \qquad
P(a\le X\le b)=\int_a^b f(x)\,dx
\]

\subsection{Typische Klausuraufgaben}
\textbf{1. Ist $f$ eine Dichte?}

Pruefen:
\begin{enumerate}
  \item $f(x)\ge 0$
  \item $\int_{-\infty}^{\infty} f(x)\,dx=1$
\end{enumerate}

\textbf{2. Aus $f$ die Verteilungsfunktion bauen}

Immer:
\[
F(x)=\int_{-\infty}^x f(u)\,du
\]

\textbf{3. Aus $F$ Wahrscheinlichkeiten berechnen}

Z.B.:
\[
P(a<X\le b)=F(b)-F(a)
\]

\textbf{4. Gleichverteilung auf Intervall}

Einfach ueber Intervalllaengen.

\subsection{Typische Fallen}
\begin{itemize}
  \item Dichtewert $f(x)$ als Wahrscheinlichkeit $P(X=x)$ missverstanden.
  \item Vergessen, dass bei stetigen Variablen $P(X=x)=0$ gilt.
  \item Bei Gleichverteilung auf $[a,b]$ falsche Hoehe der Dichte genommen.
  \item Bei Exponentialverteilung Vorzeichen im Exponenten vergeigt.
  \item Verteilungsfunktion muss monoton sein und zwischen $0$ und $1$ liegen.
\end{itemize}

\subsection{Mini-Merkliste}
\begin{itemize}
  \item $\displaystyle F(x)=P(X\le x)$
  \item jede Verteilungsfunktion ist monoton wachsend, rechtsstetig, geht von $0$ nach $1$
  \item $\displaystyle P(X=x)=F(x)-F(x^-)$
  \item stetig mit Dichte: $\displaystyle F(x)=\int_{-\infty}^x f(u)\,du$
  \item $\displaystyle P(a<X\le b)=\int_a^b f(u)\,du$
  \item bei stetigen Variablen: $\displaystyle P(X=x)=0$
  \item Gleichverteilung auf $[a,b]$: Dichte $1/(b-a)$
  \item Exponentialverteilung: $\displaystyle f(x)=\lambda e^{-\lambda x}$ fuer $x\ge 0$
\end{itemize}

\section{Kapitel 2.3: Die Normalverteilung}

\subsection{Warum wichtig?}
Die Normalverteilung taucht dauernd auf:
\begin{itemize}
  \item Messfehler
  \item technische Fertigungsmaße
  \item Summen vieler kleiner Einfluesse
  \item Approximation anderer Verteilungen
\end{itemize}

\subsection{Standardnormalverteilung}
Die Standardnormalverteilung ist
\[
\mathcal N(0,1)
\]

Dichte:
\[
\varphi(x)=\frac{1}{\sqrt{2\pi}}e^{-x^2/2}
\]

Verteilungsfunktion:
\[
\Phi(x)=P(X\le x)=\int_{-\infty}^x \varphi(u)\,du
\]

\textbf{Wichtig:}
$\Phi$ hat keine elementare Stammfunktion, deshalb arbeitet man mit Tabellen oder Rechner.

\subsection{Allgemeine Normalverteilung}
\[
X\sim \mathcal N(\mu,\sigma^2)
\qquad (\sigma>0)
\]

Dichte:
\[
f(x)=\frac{1}{\sqrt{2\pi}\sigma}\exp\left(-\frac{(x-\mu)^2}{2\sigma^2}\right)
\]

\textbf{Parameterbedeutung}
\begin{itemize}
  \item $\mu$ = Mittelwert / Lagezentrum
  \item $\sigma$ = Standardabweichung / Streuung
\end{itemize}

\subsection{Standardisierung}
Der wichtigste Schritt bei fast jeder Aufgabe:
\[
Z=\frac{X-\mu}{\sigma}
\]

Dann gilt:
\[
Z\sim \mathcal N(0,1)
\]

\textbf{Klausurregel:}
Bei Normalverteilungsaufgaben fast immer zuerst standardisieren und dann auf $\Phi$ gehen.

\subsection{Wahrscheinlichkeiten berechnen}
Wenn $X\sim \mathcal N(\mu,\sigma^2)$, dann:
\[
P(X\le x)=\Phi\left(\frac{x-\mu}{\sigma}\right)
\]

\[
P(a\le X\le b)=\Phi\left(\frac{b-\mu}{\sigma}\right)-\Phi\left(\frac{a-\mu}{\sigma}\right)
\]

\[
P(X>a)=1-\Phi\left(\frac{a-\mu}{\sigma}\right)
\]

\subsection{Symmetrie der Standardnormalverteilung}
Wegen Achsensymmetrie gilt:
\[
\Phi(-x)=1-\Phi(x)
\]

\textbf{Sehr nuetzlich:}
Tabellen enthalten oft nur Werte fuer $x\ge 0$.

\subsection{Typische Standardaufgaben}
\textbf{1. Wahrscheinlichkeit in Intervall}

Rezept:
\begin{enumerate}
  \item standardisieren
  \item auf $\Phi$ zurueckfuehren
  \item tabellieren / ablesen / rechnen
\end{enumerate}

\textbf{2. Wahrscheinlichkeit ausserhalb von Toleranzen}

Fast immer ueber Gegenereignis:
\[
P(|X-\mu|>c)=1-P(\mu-c\le X\le \mu+c)
\]

\textbf{3. Gesuchten Grenzwert $c$ bestimmen}

Wenn z.B.
\[
P(X\ge c)=0.3
\]
dann:
\[
1-\Phi\left(\frac{c-\mu}{\sigma}\right)=0.3
\]

also
\[
\Phi\left(\frac{c-\mu}{\sigma}\right)=0.7
\]
und dann aus Tabelle / inverse $\Phi$ das Quantil holen.

\subsection{$3\sigma$-Regel}
Fuer normalverteilte Variablen gilt naeherungsweise:
\[
P(\mu-3\sigma\le X\le \mu+3\sigma)\approx 0.997
\]

\textbf{Merke:}
Fast alles liegt innerhalb von drei Standardabweichungen.

\subsection{Linearkombinationen normalverteilter Variablen}
Sind $X_1,\dots,X_n$ unabhaengig und normalverteilt, dann ist auch
\[
X=a_1X_1+\dots+a_nX_n
\]
normalverteilt.

Genauer:
\[
X\sim \mathcal N(\mu,\sigma^2)
\]
mit
\[
\mu=a_1\mu_1+\dots+a_n\mu_n
\]

\[
\sigma^2=a_1^2\sigma_1^2+\dots+a_n^2\sigma_n^2
\]

\textbf{Spezialfall Summe}
Falls $X_1,X_2$ unabhaengig:
\[
X_1+X_2 \sim \mathcal N(\mu_1+\mu_2,\sigma_1^2+\sigma_2^2)
\]

\textbf{Spezialfall Differenz}
\[
X_1-X_2 \sim \mathcal N(\mu_1-\mu_2,\sigma_1^2+\sigma_2^2)
\]

\subsection{Typische Beispiele}
\textbf{1. Toleranzaufgabe}
Wenn
\[
X\sim \mathcal N(152,4)
\]
und Toleranz $\pm 5$, dann:
\[
P(147\le X\le 157)
=P\left(-2.5\le \frac{X-152}{2}\le 2.5\right)
=\Phi(2.5)-\Phi(-2.5)
\]

\textbf{2. Gewicht zweier Komponenten}
Wenn
\[
X_1\sim \mathcal N(50,1),\qquad X_2\sim \mathcal N(20,0.5)
\]
unabhaengig, dann:
\[
X_1+X_2 \sim \mathcal N(70,1.5)
\]

\subsection{Typische Fallen}
\begin{itemize}
  \item $\sigma$ und $\sigma^2$ verwechselt.
  \item Beim Standardisieren durch $\sigma^2$ statt durch $\sigma$ geteilt.
  \item Symmetrieformel falsch herum benutzt.
  \item Bei Summen normalverteilter Variablen die Varianzen falsch addiert.
  \item Toleranzaufgaben nicht ueber Gegenereignis gerechnet.
\end{itemize}

\subsection{Mini-Merkliste}
\begin{itemize}
  \item Standardnormal: $\displaystyle \varphi(x)=\frac{1}{\sqrt{2\pi}}e^{-x^2/2}$
  \item $\displaystyle X\sim \mathcal N(\mu,\sigma^2)\Rightarrow Z=\frac{X-\mu}{\sigma}\sim \mathcal N(0,1)$
  \item $\displaystyle P(X\le x)=\Phi\left(\frac{x-\mu}{\sigma}\right)$
  \item $\displaystyle P(a\le X\le b)=\Phi\left(\frac{b-\mu}{\sigma}\right)-\Phi\left(\frac{a-\mu}{\sigma}\right)$
  \item $\displaystyle \Phi(-x)=1-\Phi(x)$
  \item Summe unabhaengiger Normalvariablen bleibt normalverteilt
\end{itemize}

\section{Kapitel 2.4: Erwartungswert und Varianz bei stetigen Zufallsvariablen}

\subsection{Erwartungswert}
Hat $X$ eine Dichte $f$, dann gilt
\[
E(X)=\int_{-\infty}^{\infty} x\,f(x)\,dx
\]
falls das Integral existiert.

\textbf{Merke:}
Das ist das stetige Analogon zu
\[
E(X)=\sum_x x\cdot P(X=x).
\]

\subsection{Rechnen mit Funktionen von $X$}
Ist $Y=g(X)$, dann gilt
\[
E(g(X))=\int_{-\infty}^{\infty} g(x)\,f_X(x)\,dx
\]
falls das Integral existiert.

\textbf{Extrem wichtig:}
Du brauchst dafuer meist \textbf{nicht} erst die Dichte von $Y=g(X)$.

\subsection{Varianz}
Mit $\mu=E(X)$:
\[
D^2(X)=\operatorname{Var}(X)=E\bigl((X-\mu)^2\bigr)
\]

also bei Dichte $f$:
\[
D^2(X)=\int_{-\infty}^{\infty} (x-\mu)^2 f(x)\,dx
\]

und
\[
\sigma=\sqrt{D^2(X)}.
\]

\subsection{Wichtige Umformung}
Wie im diskreten Fall:
\[
D^2(X)=E(X^2)-\mu^2
\]
mit
\[
E(X^2)=\int_{-\infty}^{\infty} x^2 f(x)\,dx.
\]

\textbf{Merke:}
Fuer Rechnungen ist oft
\[
\operatorname{Var}(X)=E(X^2)-E(X)^2
\]
am schnellsten.

\subsection{Rechenregeln}
Die bekannten Regeln aus den diskreten Kapiteln bleiben gueltig:
\begin{itemize}
  \item $E(aX+b)=aE(X)+b$
  \item $D^2(aX+b)=a^2D^2(X)$
  \item $E(X+Y)=E(X)+E(Y)$
  \item bei Unabhaengigkeit: $E(XY)=E(X)E(Y)$
  \item $\operatorname{Var}(X+Y)=\operatorname{Var}(X)+\operatorname{Var}(Y)$ bei Unabhaengigkeit
\end{itemize}

\subsection{Standardwerte}
\textbf{Stetige Gleichverteilung auf $[a,b]$}
\[
E(X)=\frac{a+b}{2},\qquad D^2(X)=\frac{(b-a)^2}{12}
\]

\textbf{Exponentialverteilung mit Parameter $\lambda$}
\[
E(X)=\frac1\lambda,\qquad D^2(X)=\frac1{\lambda^2}
\]

\textbf{Normalverteilung $X\sim \mathcal N(\mu,\sigma^2)$}
\[
E(X)=\mu,\qquad D^2(X)=\sigma^2
\]

\subsection{Typische Klausuraufgaben}
\textbf{1. Erwartungswert aus Dichte berechnen}
\begin{enumerate}
  \item Dichte $f(x)$ sauber mit Intervall hinschreiben
  \item Integral $\int x f(x)\,dx$ ueber dem Traeger bilden
  \item sauber auswerten
\end{enumerate}

\textbf{2. Varianz aus Dichte berechnen}
\begin{enumerate}
  \item erst $\mu=E(X)$
  \item dann entweder $\int (x-\mu)^2f(x)\,dx$ oder $E(X^2)-\mu^2$
  \item am Ende $\sigma=\sqrt{\operatorname{Var}(X)}$
\end{enumerate}

\textbf{3. Erwartungswert von $g(X)$}
\[
E(g(X))=\int g(x)f_X(x)\,dx
\]
direkt benutzen.

\subsection{Typische Fallen}
\begin{itemize}
  \item ueber ganz $\mathbb R$ integriert, obwohl die Dichte nur auf einem Intervall ungleich $0$ ist.
  \item $E(g(X))$ berechnet, aber vorher unnoetig die Verteilung von $g(X)$ gesucht.
  \item Varianz mit Standardabweichung verwechselt.
  \item Dichte nicht erst auf Normierung geprueft, wenn eine Konstante $c$ vorkommt.
\end{itemize}

\subsection{Mini-Merkliste}
\begin{itemize}
  \item $\displaystyle E(X)=\int x f(x)\,dx$
  \item $\displaystyle E(g(X))=\int g(x)f_X(x)\,dx$
  \item $\displaystyle \operatorname{Var}(X)=\int (x-\mu)^2f(x)\,dx$
  \item $\displaystyle \operatorname{Var}(X)=E(X^2)-E(X)^2$
  \item $\displaystyle U[a,b]: E=\frac{a+b}{2},\ \operatorname{Var}=\frac{(b-a)^2}{12}$
  \item $\displaystyle \operatorname{Exp}(\lambda): E=\frac1\lambda,\ \operatorname{Var}=\frac1{\lambda^2}$
\end{itemize}

\section{Kapitel 2.5: Grenzwertsaetze}

\subsection{de Moivre-Laplace: Binomial $\approx$ Normal}
Wenn
\[
X\sim \operatorname{Bin}(n,p)
\]
und $n$ gross ist, dann kann man $X$ durch eine Normalverteilung ann\"ahern.

\textbf{Parameter der passenden Normalverteilung}
\[
\mu=np,\qquad \sigma^2=npq
\]

mit
\[
q=1-p
\]

\subsection{Standardisierte Form}
Man betrachtet:
\[
X^*=\frac{X-np}{\sqrt{npq}}
\]

Dann gilt fuer grosses $n$:
\[
X^* \approx \mathcal N(0,1)
\]

\textbf{Also praktisch:}
\[
P(X\le x)\approx \Phi\left(\frac{x-np}{\sqrt{npq}}\right)
\]

\subsection{Stetigkeitskorrektur}
Weil die Binomialverteilung diskret und die Normalverteilung stetig ist, verbessert man die Approximation oft durch $\pm 0.5$.

\textbf{Wichtigste Formeln}
\[
P(X\le k)\approx \Phi\left(\frac{k+0.5-np}{\sqrt{npq}}\right)
\]

\[
P(X\ge k)\approx 1-\Phi\left(\frac{k-0.5-np}{\sqrt{npq}}\right)
\]

\[
P(k\le X\le l)\approx
\Phi\left(\frac{l+0.5-np}{\sqrt{npq}}\right)
-\Phi\left(\frac{k-0.5-np}{\sqrt{npq}}\right)
\]

\textbf{Merke:}
\begin{itemize}
  \item obere Grenze bekommt meist $+0.5$
  \item untere Grenze bekommt meist $-0.5$
\end{itemize}

\subsection{Wann ist die Approximation brauchbar?}
Faustregel aus dem Blatt:
\[
npq\ge 9
\]

\textbf{Merke:}
Wenn die Varianz zu klein ist, ist die Normalapproximation oft murks.

\subsection{Typische Klausuraufgaben}
\textbf{1. Binomialwahrscheinlichkeit approximieren}

Rezept:
\begin{enumerate}
  \item $\mu=np$, $\sigma=\sqrt{npq}$
  \item passende Stetigkeitskorrektur setzen
  \item standardisieren
  \item mit $\Phi$ ausrechnen
\end{enumerate}

\textbf{2. Mindestens / hoechstens / zwischen}

Sauber in Intervallform bringen und dann die richtige $\pm 0.5$-Korrektur setzen.

\subsection{Typische Fallen}
\begin{itemize}
  \item Stetigkeitskorrektur vergessen.
  \item $\sigma=\sqrt{npq}$ mit $\sigma^2=npq$ verwechselt.
  \item Approximation benutzt, obwohl $npq$ zu klein ist.
  \item Bei ``mindestens'' und ``hoechstens'' die falsche Grenze korrigiert.
\end{itemize}

\subsection{Mini-Merkliste}
\begin{itemize}
  \item $\displaystyle X\sim \operatorname{Bin}(n,p)\approx \mathcal N(np,npq)$
  \item $\displaystyle Z=\frac{X-np}{\sqrt{npq}}$
  \item mit Stetigkeitskorrektur arbeiten
  \item Faustregel: $\displaystyle npq\ge 9$
\end{itemize}

\section{Kapitel 3.1: Statistik-Ueberblick}

\subsection{Wichtigste Grundbegriffe}
\textbf{Grundgesamtheit}
\begin{itemize}
  \item Menge aller Objekte, fuer die man sich interessiert
\end{itemize}

\textbf{Statistische Einheit}
\begin{itemize}
  \item ein einzelnes Element der Grundgesamtheit
\end{itemize}

\textbf{Merkmal}
\begin{itemize}
  \item eine beobachtete Eigenschaft einer statistischen Einheit
\end{itemize}

\textbf{Auspraegung / Wert}
\begin{itemize}
  \item konkreter Wert, den das Merkmal annimmt
\end{itemize}

\textbf{Stichprobe}
\begin{itemize}
  \item die tatsaechlich untersuchten Einheiten
\end{itemize}

\subsection{Uni-, bi-, multivariat}
\begin{itemize}
  \item \textbf{univariat}: ein Merkmal
  \item \textbf{bivariat}: zwei Merkmale
  \item \textbf{multivariat}: mehrere Merkmale
\end{itemize}

\subsection{Skalenniveaus}
\textbf{Nominal skaliert}
\begin{itemize}
  \item nur gleich / verschieden
  \item Beispiel: Geschlecht
\end{itemize}

\textbf{Ordinal skaliert}
\begin{itemize}
  \item Reihenfolge moeglich
  \item Abstaende nicht sinnvoll interpretierbar
  \item Beispiel: Schulnoten
\end{itemize}

\textbf{Metrisch skaliert}
\begin{itemize}
  \item Zahlenwerte mit sinnvollen Differenzen
  \item Beispiel: Gewicht, Durchmesser, Alter
\end{itemize}

\subsection{Diskret vs. kontinuierlich}
\textbf{Diskret}
\begin{itemize}
  \item endlich oder abz\"ahlbar viele moegliche Werte
\end{itemize}

\textbf{Kontinuierlich / stetig}
\begin{itemize}
  \item Werte liegen in einem lueckenlosen Bereich / Intervall
\end{itemize}

\subsection{Arten der Statistik}
\textbf{Beschreibende Statistik}
\begin{itemize}
  \item Daten aufbereiten, darstellen, zusammenfassen
\end{itemize}

\textbf{Schliessende Statistik}
\begin{itemize}
  \item aus Stichproben auf die Grundgesamtheit schliessen
  \item dabei bleibt ein Irrtumsrisiko
\end{itemize}

\subsection{Bezug zur Wahrscheinlichkeitstheorie}
Bei Zufallsstichproben:
\begin{itemize}
  \item Grundgesamtheit $\leftrightarrow \Omega$
  \item statistische Einheit $\leftrightarrow \omega\in\Omega$
  \item Merkmal $\leftrightarrow$ Zufallsvariable $X$
  \item Stichprobe $\leftrightarrow (\omega_1,\dots,\omega_n)\in \Omega^n$
\end{itemize}

\subsection{Mini-Merkliste}
\begin{itemize}
  \item Grundgesamtheit = alle interessierenden Objekte
  \item Stichprobe = untersuchte Auswahl
  \item nominal / ordinal / metrisch unterscheiden
  \item diskret vs. stetig unterscheiden
  \item beschreibend = Daten darstellen
  \item schliessend = von Stichprobe auf Grundgesamtheit schliessen
\end{itemize}

\section{Kapitel 3.2: Univariate beschreibende Statistik}

\subsection{Grundnotation}
Gegeben sei eine Stichprobe
\[
x_1,\dots,x_n
\]
mit Stichprobenumfang $n$.

Die verschiedenen beobachteten Werte seien
\[
a_1< a_2 < \dots < a_k.
\]

\textbf{Notation:}
\begin{itemize}
  \item $n_i$ = absolute Haeufigkeit von $a_i$
  \item $r_i=\frac{n_i}{n}$ = relative Haeufigkeit von $a_i$
  \item stets: $\sum_{i=1}^k n_i=n$ und $\sum_{i=1}^k r_i=1$
\end{itemize}

\subsection{Relative Haeufigkeitsfunktion}
\[
\hat f(x)=
\begin{cases}
r_i,& x=a_i\\
0,& \text{sonst}
\end{cases}
\]

\textbf{Bedeutung:}
\begin{itemize}
  \item $\hat f(x)$ = Anteil der Stichprobendaten mit Wert $x$
  \item diskretes Gegenstueck zu $P(X=x)$
\end{itemize}

\subsection{Empirische Verteilungsfunktion}
\[
\hat F(x)=\frac{1}{n}\#\{j: x_j\le x\}
\]

aequivalent:
\[
\hat F(x)=
\begin{cases}
0,& x<a_1\\
r_1+\dots+r_i,& a_i\le x<a_{i+1}\\
1,& x\ge a_k
\end{cases}
\]

\textbf{Merke:}
\begin{itemize}
  \item $\hat F$ ist eine Treppenfunktion.
  \item Spruenge liegen bei $a_i$.
  \item Sprunghoehe an $a_i$ ist $r_i=\hat f(a_i)$.
\end{itemize}

\subsection{Bezug zur Wahrscheinlichkeit}
Bei einer Zufallsstichprobe gilt fuer grosses $n$ heuristisch:
\[
\hat f(x)\approx P(X=x),\qquad \hat F(x)\approx F(x)=P(X\le x).
\]

\textbf{Merke:}
Relative Haeufigkeit ist die empirische Version von Wahrscheinlichkeit.

\subsection{Datenfriedhof und Klasseneinteilung}
Bei vielen verschiedenen Werten ist das Stabdiagramm oft unbrauchbar. Dann in Klassen
\[
I_1,\dots,I_k
\]
einteilen.

Fuer Klasse $I_i$ mit Breite $\Delta x_i$:
\begin{itemize}
  \item $n_i$ = Klassenhaeufigkeit
  \item $r_i=\frac{n_i}{n}$ = relative Klassenhaeufigkeit
  \item $f_i=\frac{r_i}{\Delta x_i}$ = Klassenhaeufigkeitsdichte
  \item $F_i=r_1+\dots+r_i=\sum_{j=1}^i f_j\Delta x_j$
\end{itemize}

\textbf{Faustregel aus dem Blatt:}
\[
k\approx 2\ln(n)\qquad (n\ge 50)
\]

\subsection{Histogramm}
Im Histogramm ist ueber Klasse $I_i$ ein Balken mit
\begin{itemize}
  \item Breite $\Delta x_i$
  \item Hoehe $f_i=\frac{r_i}{\Delta x_i}$
\end{itemize}
einzuzeichnen.

Dann gilt:
\[
\text{Flaeche des Balkens}=f_i\Delta x_i=r_i.
\]

\textbf{Extrem wichtig:}
\begin{itemize}
  \item Nicht $r_i$ als Hoehe nehmen, sondern $f_i$.
  \item Im Histogramm repraesentiert die \textbf{Flaeche} die relative Haeufigkeit.
  \item Summe aller Balkenflaechen = $1$.
\end{itemize}

\subsection{Lagemasse}
\textbf{Arithmetisches Mittel}
\[
\bar x=\frac{1}{n}\sum_{j=1}^n x_j=\sum_{i=1}^k r_i a_i
\]

\textbf{Median}

Zuerst Daten sortieren:
\[
x_{(1)}\le x_{(2)}\le \dots \le x_{(n)}.
\]

Dann:
\[
\tilde x=
\begin{cases}
x_{(\frac{n+1}{2})},& n\ \text{ungerade}\\[1mm]
\frac12\bigl(x_{(\frac n2)}+x_{(\frac n2+1)}\bigr),& n\ \text{gerade}
\end{cases}
\]

\textbf{Merke:}
\begin{itemize}
  \item $\bar x$ reagiert empfindlich auf Ausreisser.
  \item Median ist robuster gegen Ausreisser.
\end{itemize}

\subsection{Streuungsmasse}
\textbf{Spannweite}
\[
R=x_{\max}-x_{\min}
\]

\textbf{Empirische Varianz}
\[
s^2=\frac{1}{n-1}\sum_{j=1}^n (x_j-\bar x)^2
\]

\textbf{Empirische Standardabweichung}
\[
s=\sqrt{s^2}
\]

\textbf{Rechenhilfe}
\[
\sum_{j=1}^n (x_j-\bar x)^2=\sum_{j=1}^n x_j^2-n\bar x^2
\]

also
\[
s^2=\frac{1}{n-1}\left(\sum_{j=1}^n x_j^2-n\bar x^2\right).
\]

\textbf{Typischer Fall:}
Wenn in der Aufgabe $x_1,\dots,x_n$ direkt gegeben sind, erst $\bar x$, dann $\sum x_j^2$, dann $s^2$.

\subsection{Typische Klausuraufgaben}
\textbf{1. Haeufigkeitstabelle bauen}
\begin{enumerate}
  \item verschiedene Werte $a_i$ sortieren
  \item $n_i$ zaehlen
  \item $r_i=n_i/n$
  \item kumulieren fuer $\hat F$
\end{enumerate}

\textbf{2. Histogramm zeichnen}
\begin{enumerate}
  \item Klassen festlegen
  \item $n_i$, $r_i$ bestimmen
  \item Hoehe $f_i=r_i/\Delta x_i$
  \item Flaechen checken
\end{enumerate}

\textbf{3. Lage- und Streuungsmasse}
\begin{enumerate}
  \item $\bar x$ berechnen
  \item Median aus sortierten Daten
  \item $R$, $s^2$, $s$
\end{enumerate}

\subsection{Typische Fallen}
\begin{itemize}
  \item Histogrammhoehe mit relativer Haeufigkeit statt Dichte verwechselt.
  \item Bei geradem $n$ den Median falsch genommen.
  \item In $s^2$ versehentlich durch $n$ statt durch $n-1$ geteilt.
  \item Unsortierte Daten fuer den Median benutzt.
\end{itemize}

\subsection{Mini-Merkliste}
\begin{itemize}
  \item $r_i=\frac{n_i}{n}$, \ $\hat f(a_i)=r_i$
  \item $\hat F(x)=\frac1n\#\{j:x_j\le x\}$
  \item Histogramm: Hoehe $f_i=\frac{r_i}{\Delta x_i}$, Flaeche $=r_i$
  \item $\bar x=\frac1n\sum x_j$
  \item $s^2=\frac{1}{n-1}\sum (x_j-\bar x)^2$
  \item $s=\sqrt{s^2}$
\end{itemize}

\section{Kapitel 4.1: Parameterschaetzung}

\subsection{Worum geht es?}
Die Verteilung einer Zufallsvariable haengt von unbekannten Parametern ab, z.B.
\begin{itemize}
  \item Bernoulli / Binomial: $p$
  \item Normalverteilung: $\mu$, $\sigma^2$
\end{itemize}

Ziel:
Diese unbekannten Parameter aus einer Stichprobe schaetzen.

\subsection{Standardmodell fuer Anteilsprobleme}
Bei Ja/Nein-Eigenschaft:
\[
X(\omega)=
\begin{cases}
1,& \text{Eigenschaft liegt vor}\\
0,& \text{sonst}
\end{cases}
\]

Dann gilt:
\[
X\sim \operatorname{Bern}(p),\qquad E(X)=p,\qquad D^2(X)=pq,\quad q=1-p.
\]

Bei einer Zufallsstichprobe:
\[
X_1,\dots,X_n \text{ iid wie } X.
\]

\subsection{Relative Haeufigkeit als Schluesselgroesse}
Die Stichprobenmittelzahl
\[
\bar X=\frac1n\sum_{i=1}^n X_i
\]
ist bei Bernoulli-Daten genau die relative Haeufigkeit der interessierenden Eigenschaft.

Aus frueheren Kapiteln:
\[
E(\bar X)=p,\qquad D^2(\bar X)=\frac{pq}{n}.
\]

\subsection{Idee von Punkt- und Intervallschaetzung}
\textbf{Punktschaetzung:}
Ein einzelner Wert, z.B.
\[
\hat p=\bar X.
\]

\textbf{Intervallschaetzung:}
Ein zufaelliges Intervall um den Schaetzwert, z.B.
\[
[\bar X-c,\bar X+c].
\]

\subsection{Warum funktioniert das?}
Mit Tschebyscheff:
\[
P(|\bar X-p|>c)\le \frac{pq}{nc^2}\le \frac{1}{4nc^2}.
\]

Also
\[
P(|\bar X-p|\le c)\ge 1-\frac{pq}{nc^2}\ge 1-\frac{1}{4nc^2}.
\]

\textbf{Merke:}
\begin{itemize}
  \item groesseres $n$ $\Rightarrow$ genauere Schaetzung
  \item kleineres $c$ $\Rightarrow$ engeres, aber schwerer erreichbares Intervall
  \item $\gamma$ = Vertrauensniveau, oft $0.95$ oder $0.99$
\end{itemize}

\subsection{Typische Notation}
\begin{itemize}
  \item wahrer unbekannter Parameter: $p$, $\mu$, $\sigma^2$, allgemein $\vartheta$
  \item Schaetzwert: $\hat p$, $\hat\mu$, $\widehat{\sigma^2}$, allgemein $\hat\vartheta$
  \item Stichprobe: $X_1,\dots,X_n$
  \item beobachtete Werte: $x_1,\dots,x_n$
\end{itemize}

\subsection{Typische Klausurideen}
\textbf{1. Anteil modellieren}
\begin{itemize}
  \item Indikatorvariable $X\in\{0,1\}$ definieren
  \item dann Bernoulli-Modell benutzen
\end{itemize}

\textbf{2. Schaetzidee begruenden}
\begin{itemize}
  \item $\bar X$ liegt nahe bei $p$, weil $E(\bar X)=p$ und $D^2(\bar X)=pq/n\to 0$
\end{itemize}

\textbf{3. noetigen Stichprobenumfang abschaetzen}
\[
1-\frac{1}{4nc^2}\ge \gamma
\]
nach $n$ aufloesen.

\subsection{Mini-Merkliste}
\begin{itemize}
  \item unbekannter Parameter soll aus Stichprobe geschaetzt werden
  \item bei Anteilen: $X\sim \operatorname{Bern}(p)$
  \item $\bar X=\frac1n\sum X_i$ ist relative Haeufigkeit
  \item $E(\bar X)=p$, \ $D^2(\bar X)=pq/n$
  \item Punktschaetzung = einzelner Wert, Intervallschaetzung = Bereich
\end{itemize}

\section{Kapitel 4.2: Punktschaetzung}

\subsection{Allgemeine Form}
Eine Schaetzfunktion ist eine Funktion der Stichprobe:
\[
\hat\vartheta=t_n(x_1,\dots,x_n).
\]

Als Zufallsvariable geschrieben:
\[
T_n=t_n(X_1,\dots,X_n).
\]

\textbf{Merke:}
\begin{itemize}
  \item $\vartheta$ = wahrer Parameter
  \item $\hat\vartheta$ = konkreter Schaetzwert aus Daten
  \item $T_n$ = zugehoeriger Schaetzer als Zufallsvariable
\end{itemize}

\subsection{Wichtige Standardschaetzer}
\textbf{Unbekannter Mittelwert $\mu$}
\[
\hat\mu=\bar X=\frac1n\sum_{i=1}^n X_i
\]

Eigenschaften:
\[
E(\bar X)=\mu,\qquad D^2(\bar X)=\frac{\sigma^2}{n}.
\]

\textbf{Unbekannte Varianz $\sigma^2$}
\[
S^2=\frac{1}{n-1}\sum_{i=1}^n (X_i-\bar X)^2
\]

und fuer beobachtete Daten:
\[
s^2=\frac{1}{n-1}\sum_{i=1}^n (x_i-\bar x)^2.
\]

\textbf{Extrem wichtig:}
\[
E(S^2)=\sigma^2.
\]

Darum nimmt man bei Schaetzung der Varianz
\[
\frac{1}{n-1}\sum (X_i-\bar X)^2
\]
und nicht die Version mit $\frac1n$.

\subsection{Erwartungstreue}
Definition:
\[
E(T_n)=\vartheta
\]

\textbf{Bedeutung:}
Der Schaetzer trifft den wahren Parameter im Mittel richtig.

\textbf{Bias / Verzerrung:}
\[
\operatorname{Bias}(T_n)=E(T_n)-\vartheta.
\]

\textbf{Merke:}
\begin{itemize}
  \item erwartungstreu $\Leftrightarrow \operatorname{Bias}(T_n)=0$
  \item $\bar X$ ist erwartungstreu fuer $\mu$
  \item $S^2$ ist erwartungstreu fuer $\sigma^2$
\end{itemize}

\subsection{Konsistenz}
Definition:
\[
\forall \varepsilon>0:\quad P(|T_n-\vartheta|\ge \varepsilon)\to 0 \qquad (n\to\infty)
\]

\textbf{Bedeutung:}
Mit wachsendem $n$ landet der Schaetzer immer wahrscheinlicher nahe am wahren Wert.

\textbf{Merke:}
Konsistenz ist das ``langfristig richtig''-Kriterium.

\subsection{Effizienz}
Sind $T_n$ und $T_n'$ beide erwartungstreu, dann ist der Schaetzer mit kleinerer Varianz besser.

\[
D^2(T_n)<D^2(T_n')
\]
bedeutet:
$T_n$ streut weniger um den wahren Wert.

\textbf{Definition:}
Effizient = erwartungstreu und unter den erwartungstreuen Schaetzern minimale Varianz.

\subsection{Typische Klausuraufgaben}
\textbf{1. Erwartungstreue pruefen}
\begin{enumerate}
  \item Schaetzer $T_n$ hinschreiben
  \item $E(T_n)$ mit Linearitaet ausrechnen
  \item mit $\vartheta$ vergleichen
\end{enumerate}

\textbf{2. Besseren Schaetzer waehlen}
\begin{enumerate}
  \item erst Erwartungstreue pruefen
  \item dann Varianzen vergleichen
  \item kleinere Varianz gewinnt
\end{enumerate}

\textbf{3. Varianzschaetzer erkennen}
\begin{itemize}
  \item fuer Erwartungstreue der Varianz steht fast immer der Faktor $\frac{1}{n-1}$ da
\end{itemize}

\subsection{Rechenmuster}
\textbf{Linearitaet des Erwartungswerts}
\[
E\left(\sum a_iX_i+c\right)=\sum a_iE(X_i)+c
\]

\textbf{Varianz bei Unabhaengigkeit}
\[
D^2\left(\sum a_iX_i\right)=\sum a_i^2D^2(X_i)
\]

damit pruefst du fast alle Aufgaben zu vorgegebenen Schaetzern.

\subsection{Typische Fallen}
\begin{itemize}
  \item Schaetzwert $\hat\vartheta$ und Schaetzer $T_n$ vermischt.
  \item Nur Erwartungstreue geprueft, aber nicht die Streuung.
  \item Bei Varianzschaetzung den Faktor $n-1$ vergessen.
  \item Konsistenz mit Erwartungstreue verwechselt.
\end{itemize}

\subsection{Mini-Merkliste}
\begin{itemize}
  \item $\hat\vartheta=t_n(x_1,\dots,x_n)$
  \item $T_n=t_n(X_1,\dots,X_n)$
  \item erwartungstreu: $E(T_n)=\vartheta$
  \item konsistent: $P(|T_n-\vartheta|\ge\varepsilon)\to 0$
  \item effizient: unter erwartungstreuen Schaetzern kleinste Varianz
  \item $\hat\mu=\bar X$, \ $\widehat{\sigma^2}=S^2=\frac{1}{n-1}\sum (X_i-\bar X)^2$
\end{itemize}

\section{Kapitel 4.4: Hypothesentest}

\subsection{Grundidee}
Ein Hypothesentest ist eine Entscheidungsregel auf Basis einer Stichprobe:
\begin{itemize}
  \item $H_0$ = Nullhypothese
  \item $H_1$ = Alternative
\end{itemize}

Am Ende wird entweder
\begin{itemize}
  \item $H_0$ beibehalten oder
  \item $H_0$ verworfen
\end{itemize}

\textbf{Merke:}
``$H_0$ verwerfen'' bedeutet gleichzeitig ``$H_1$ annehmen''.

\subsection{Testfunktion und kritischer Bereich}
Ein Test basiert auf einer Prueffunktion / Testfunktion
\[
T_n=t_n(X_1,\dots,X_n).
\]

Der Wertebereich wird zerlegt in:
\begin{itemize}
  \item $O$ = OK-Bereich / Annahmebereich fuer $H_0$
  \item $K$ = kritischer Bereich / Ablehnungsbereich fuer $H_0$
\end{itemize}

Entscheidungsregel:
\[
T_n\in O \Rightarrow H_0 \text{ beibehalten},\qquad
T_n\in K \Rightarrow H_0 \text{ verwerfen}.
\]

\subsection{Fehler 1. und 2. Art}
\textbf{Fehler 1. Art}
\[
H_0 \text{ ist wahr, wird aber verworfen}
\]
Wahrscheinlichkeit hoechstens
\[
\alpha.
\]

\textbf{Fehler 2. Art}
\[
H_1 \text{ ist wahr, aber } H_0 \text{ wird beibehalten}
\]
oft mit Wahrscheinlichkeit $\beta$ abgeschaetzt.

\textbf{Extrem wichtig:}
\begin{itemize}
  \item $\alpha$ wird vorgegeben.
  \item $\beta$ ist meist nicht gleichzeitig frei kontrollierbar.
  \item $1-\alpha$ ist \textbf{nicht} die Wahrscheinlichkeit, dass $H_1$ wahr ist.
\end{itemize}

\subsection{Rollen von $H_0$ und $H_1$}
Die Rollen sind \textbf{nicht} austauschbar.

\textbf{Merke:}
In $H_0$ steckt die Aussage, die man nur bei genuegend starkem Widerspruch verwirft.

Typische Formen:
\begin{itemize}
  \item zweiseitig: $H_0:\mu=\mu_0$ gegen $H_1:\mu\ne\mu_0$
  \item rechtsseitig: $H_0:p\le p_0$ gegen $H_1:p>p_0$
  \item linksseitig: $H_0:\mu\ge \mu_0$ gegen $H_1:\mu<\mu_0$
\end{itemize}

\subsection{Guetefunktion und Testcharakteristik}
\[
G(\vartheta)=P_\vartheta(T_n\in K)
\]
= Wahrscheinlichkeit, dass der Test $H_0$ verwirft.

\[
L(\vartheta)=P_\vartheta(T_n\in O)
\]
= Wahrscheinlichkeit, dass der Test $H_0$ beibehaelt.

Es gilt:
\[
L(\vartheta)+G(\vartheta)=1.
\]

\subsection{Wichtiger Standardtest}
\textbf{Zweiseitiger Test fuer den Mittelwert einer normalverteilten Variable bei bekannter Varianz}

Voraussetzungen:
\[
X_1,\dots,X_n \text{ iid } \sim \mathcal N(\mu,\sigma^2),\qquad \sigma^2 \text{ bekannt}
\]

Hypothesen:
\[
H_0:\mu=\mu_0,\qquad H_1:\mu\ne \mu_0
\]

Testfunktion:
\[
T_n=\bar X=\frac1n\sum_{i=1}^n X_i
\]

unter $H_0$ gilt:
\[
\bar X\sim \mathcal N\left(\mu_0,\frac{\sigma^2}{n}\right).
\]

\subsection{Kritischer Bereich beim zweiseitigen $\mu$-Test}
Mit Signifikanzniveau $\alpha$:
\[
K=\left(-\infty,\mu_0-c\right)\cup \left(\mu_0+c,\infty\right)
\]

mit
\[
c=\frac{\sigma}{\sqrt n}\,u_{1-\alpha/2}
\]

wobei $u_{1-\alpha/2}$ das $(1-\alpha/2)$-Quantil der Standardnormalverteilung ist.

Also:
\[
H_0 \text{ verwerfen } \Longleftrightarrow |\bar X-\mu_0|> \frac{\sigma}{\sqrt n}u_{1-\alpha/2}.
\]

aequivalent standardisiert:
\[
Z=\frac{\bar X-\mu_0}{\sigma/\sqrt n}
\]
und
\[
H_0 \text{ verwerfen } \Longleftrightarrow |Z|>u_{1-\alpha/2}.
\]

\subsection{Wichtige Quantile}
\begin{itemize}
  \item bei $\alpha=0.05$: $u_{0.975}\approx 1.96$
  \item bei $\alpha=0.01$: $u_{0.995}\approx 2.576$
\end{itemize}

\subsection{Typische Klausuraufgaben}
\textbf{1. Ablehnen oder nicht?}
\begin{enumerate}
  \item Hypothesen hinschreiben
  \item $\bar x$ berechnen
  \item Testgroesse $Z=\frac{\bar x-\mu_0}{\sigma/\sqrt n}$ berechnen
  \item mit kritischem Wert vergleichen
\end{enumerate}

\textbf{2. Minimal noetiges $n$ finden}
Aus
\[
|\bar x-\mu_0|>\frac{\sigma}{\sqrt n}u_{1-\alpha/2}
\]
nach $n$ aufloesen:
\[
n>\left(\frac{\sigma\,u_{1-\alpha/2}}{|\bar x-\mu_0|}\right)^2.
\]

\textbf{3. Testart erkennen}
\begin{itemize}
  \item ``ungleich'' $\Rightarrow$ zweiseitig
  \item ``groesser'' $\Rightarrow$ rechtsseitig
  \item ``kleiner'' $\Rightarrow$ linksseitig
\end{itemize}

\subsection{Typische Fallen}
\begin{itemize}
  \item $H_0$ und $H_1$ vertauscht.
  \item Aus ``$H_0$ nicht verworfen'' faelschlich ``$H_0$ ist bewiesen'' gemacht.
  \item $\alpha$ als Wahrscheinlichkeit interpretiert, dass $H_0$ oder $H_1$ wahr ist.
  \item Bei zweiseitigem Test $u_{1-\alpha}$ statt $u_{1-\alpha/2}$ benutzt.
  \item $\sigma/\sqrt n$ vergessen.
\end{itemize}

\subsection{Mini-Merkliste}
\begin{itemize}
  \item Test = Entscheidung zwischen $H_0$ und $H_1$
  \item Fehler 1. Art: wahres $H_0$ verworfen, Wkt. $\le \alpha$
  \item Fehler 2. Art: falsches $H_0$ beibehalten
  \item ``nicht verworfen'' heisst nicht ``bewiesen''
  \item bei $H_0:\mu=\mu_0$ und bekannter $\sigma^2$:
  \[
  Z=\frac{\bar X-\mu_0}{\sigma/\sqrt n}
  \]
  \item zweiseitig ablehnen, wenn $\displaystyle |Z|>u_{1-\alpha/2}$
\end{itemize}

\end{document}
